% =====================================================================
%  TP2 — Infrastructure as Code avec Terraform sur Azure (ShopEasy)
%  Binôme : Louis SCARFONE & Maxence BOURRAGUE — Bloc 4 (Cloud Computing)
%  Moteur : LuaLaTeX (LuaHBTeX)
% =====================================================================
\documentclass[11pt,a4paper]{article}

% ---------- Langue & fontes ----------
\usepackage{fontspec}
\usepackage{polyglossia}
\setdefaultlanguage{french}
\setmainfont{TeX Gyre Pagella}
\setsansfont{TeX Gyre Heros}
\setmonofont{DejaVu Sans Mono}[Scale=0.80]
\usepackage{microtype}
\emergencystretch=3em  % absorbe les débordements de texte (code inline en fin de ligne)

% ---------- Respiration (look épuré) ----------
\usepackage{setspace}
\setstretch{1.08}
\setlength{\parindent}{0pt}
\setlength{\parskip}{0.7em plus 0.12em minus 0.05em}

% ---------- Géométrie ----------
\usepackage[a4paper,margin=2.2cm,headheight=16pt,headsep=14pt,footskip=20pt]{geometry}

% ---------- Couleurs (palette Azure) ----------
\usepackage{xcolor}
\definecolor{azure}{HTML}{0078D4}
\definecolor{azuredark}{HTML}{14375E}
\definecolor{azuremid}{HTML}{2E75B6}
\definecolor{azurelight}{HTML}{E8F1FB}
\definecolor{azurepale}{HTML}{F3F8FD}
\definecolor{okgreen}{HTML}{107C10}
\definecolor{okgreenbg}{HTML}{E8F5E9}
\definecolor{warnorange}{HTML}{B5520A}
\definecolor{warnbg}{HTML}{FBEEDF}
\definecolor{dangerred}{HTML}{A4262C}
\definecolor{codebg}{HTML}{F6F8FA}
\definecolor{codeframe}{HTML}{D8E1EC}
\definecolor{ink}{HTML}{1B2733}
\definecolor{rowalt}{HTML}{EBF3FC}
\color{ink}

% ---------- Divers ----------
\usepackage{graphicx}
\graphicspath{{../livrables/assets/}}
\usepackage{fontawesome5}
\usepackage{enumitem}
\setlist{leftmargin=1.4em,topsep=2pt,itemsep=2pt,parsep=0pt}
\usepackage{needspace}
\usepackage{fvextra}

% ---------- Emojis (Twemoji couleur) ----------
\usepackage{twemojis}
\newcommand{\emo}[1]{\raisebox{-0.14em}{\twemoji{#1}}}
\newcommand{\eTarget}{\emo{1f3af}}
\newcommand{\ePackage}{\emo{1f4e6}}
\newcommand{\eCheck}{\emo{2705}}
\newcommand{\eWarn}{\emo{26a0}}
\newcommand{\eBulb}{\emo{1f4a1}}
\newcommand{\eLock}{\emo{1f512}}
\newcommand{\eKey}{\emo{1f510}}
\newcommand{\eRocket}{\emo{1f680}}
\newcommand{\eChart}{\emo{1f4ca}}
\newcommand{\eMoney}{\emo{1f4b0}}
\newcommand{\eGear}{\emo{2699}}
\newcommand{\eGlobe}{\emo{1f310}}
\newcommand{\eCamera}{\emo{1f4f8}}
\newcommand{\ePin}{\emo{1f4cc}}
\newcommand{\eMag}{\emo{1f50d}}
\newcommand{\eFlag}{\emo{1f3c1}}
\newcommand{\eParty}{\emo{1f389}}
\newcommand{\eShield}{\emo{1f6e1}}
\newcommand{\eBroom}{\emo{1f9f9}}
\newcommand{\eFolder}{\emo{1f4c2}}
\newcommand{\eServer}{\emo{1f5a5}}
\newcommand{\eDB}{\emo{1f5c4}}
\newcommand{\eScale}{\emo{2696}}
\newcommand{\eRed}{\emo{1f534}}
\newcommand{\eOrange}{\emo{1f7e0}}
\newcommand{\eYellow}{\emo{1f7e1}}
\newcommand{\eGreen}{\emo{1f7e2}}
\newcommand{\eWorld}{\emo{1f30d}}
\newcommand{\arr}{\ensuremath{\rightarrow}}
\newcommand{\code}[1]{\texttt{#1}}

% ---------- Tableaux (tabularray) ----------
\usepackage{tabularray}
\UseTblrLibrary{booktabs}
\NewDocumentEnvironment{ltable}{m}{%
  \par\addvspace{4pt}\begin{tblr}{
    width=\linewidth,
    colspec={#1},
    row{1}={bg=azuredark,fg=white,font=\sffamily\bfseries\small},
    row{2-Z}={font=\small},
    row{even}={bg=rowalt},
    hline{1,Z}={1.1pt,solid,azuredark},
    hline{2}={0.5pt,solid,azuredark!45},
    colsep=7pt, rowsep=4pt,
    stretch=1.18,
  }%
}{%
  \end{tblr}\par\addvspace{4pt}%
}

% ---------- Boîtes (tcolorbox) ----------
\usepackage[most]{tcolorbox}
\tcbset{
  boxbase/.style={
    enhanced, breakable, sharp corners, boxrule=0pt, leftrule=4pt,
    left=4.5mm, right=4.5mm, top=3.6mm, bottom=3.6mm,
    before skip=12pt, after skip=12pt,
    fonttitle=\sffamily\bfseries, coltitle=white,
    attach boxed title to top left={xshift=0mm,yshift=-2.6mm},
    boxed title style={sharp corners, boxrule=0pt, left=2.5mm,right=2.5mm,top=0.9mm,bottom=0.9mm},
  },
}
\newtcolorbox{fiche}{boxbase, colback=azurelight, colframe=azure}
\newtcolorbox{notebox}[1][Point de vigilance]{boxbase, colback=warnbg, colframe=warnorange,
  colbacktitle=warnorange, title={\eWarn~#1}}
\newtcolorbox{tipbox}[1][Astuce]{boxbase, colback=azurepale, colframe=azuremid,
  colbacktitle=azuremid, title={\eBulb~#1}}
\newtcolorbox{etatbox}[1][État après l'atelier]{boxbase, colback=okgreenbg, colframe=okgreen,
  colbacktitle=okgreen, title={\eCheck~#1}}

% Bloc code / terminal
\fvset{breaklines=true, breakanywhere=true, fontsize=\footnotesize,
  breaksymbolleft=\raisebox{0.1ex}{\scriptsize\textcolor{azure}{$\hookrightarrow$}}}
\newtcolorbox{codebox}[1]{
  enhanced, breakable, sharp corners, boxrule=0.7pt,
  colback=codebg, colframe=codeframe,
  coltitle=azuredark, colbacktitle=codeframe, fonttitle=\ttfamily\bfseries\scriptsize,
  title={\faTerminal~~#1}, left=3mm, right=3mm, top=2.6mm, bottom=2.6mm,
  before skip=12pt, after skip=14pt,
}

% Capture d'écran encadrée + légende
\newcommand{\screenshot}[2]{%
  \par\addvspace{10pt}
  \begin{center}
  \begin{tcolorbox}[enhanced, sharp corners, boxrule=0.7pt, colframe=codeframe, colback=white,
      left=0.8mm, right=0.8mm, top=0.8mm, bottom=0.8mm, drop fuzzy shadow=black!18]
    \includegraphics[width=\linewidth]{#1}
  \end{tcolorbox}\par\vspace{4pt}
  {\small\itshape\color{black!58}#2}
  \end{center}\par\addvspace{12pt}
}

% Carte « chiffre-clé » pour la page de garde
\newcommand{\statcard}[2]{%
  \begin{tcolorbox}[enhanced, sharp corners, boxrule=0pt, leftrule=3pt, colframe=azure,
     colback=azurepale, width=0.305\linewidth, halign=center, valign=center, height=2.3cm,
     left=1mm, right=1mm, top=1mm, bottom=1mm, nobeforeafter]
     {\sffamily\bfseries\fontsize{22}{24}\selectfont\color{azuredark}#1}\\[1.5mm]
     {\sffamily\small\color{black!58}#2}
  \end{tcolorbox}%
}

% ---------- Sections (titlesec) ----------
\usepackage{titlesec}
\titleformat{\section}
  {\sffamily\LARGE\bfseries\color{azuredark}}
  {\tcbox[on line, sharp corners, boxrule=0pt, colback=azure, colupper=white,
     left=2.2mm,right=2.2mm,top=0.6mm,bottom=0.6mm]{\thesection}}
  {0.7em}{}[\vspace{1mm}{\color{azure}\titlerule[1.6pt]}]
\titlespacing*{\section}{0pt}{3.6ex plus 0.8ex}{2.0ex}
\titleformat{\subsection}
  {\sffamily\large\bfseries\color{azuremid}}{\thesubsection}{0.6em}{}
\titlespacing*{\subsection}{0pt}{2.8ex plus 0.5ex}{1.1ex}
\titleformat{\subsubsection}
  {\sffamily\bfseries\color{azuredark}}{}{0pt}{}
\titlespacing*{\subsubsection}{0pt}{1.9ex plus 0.3ex}{0.5ex}

% ---------- En-têtes / pieds (fancyhdr) ----------
\usepackage{fancyhdr}
\pagestyle{fancy}
\fancyhf{}
\fancyhead[L]{\sffamily\small\color{azuredark}\textbf{TP2 Terraform}~\textcolor{azure}{\textbullet}~ShopEasy}
\fancyhead[R]{\sffamily\small\color{black!55}\nouppercase{\leftmark}}
\fancyfoot[C]{\sffamily\small\color{black!55}\thepage}
\fancyfoot[R]{\sffamily\scriptsize\color{black!42}Bloc 4 · Cloud Computing}
\fancyfoot[L]{\sffamily\scriptsize\color{black!42}SCARFONE · BOURRAGUE}
\renewcommand{\headrulewidth}{0pt}
\renewcommand{\footrulewidth}{0pt}
\renewcommand{\headrule}{\vspace{-2pt}{\color{azure!70}\rule{\headwidth}{1.1pt}}}

% ---------- Liens ----------
\usepackage[hidelinks]{hyperref}
\hypersetup{colorlinks=true, linkcolor=azuredark, urlcolor=azuremid,
  pdftitle={TP2 Terraform - ShopEasy}, pdfauthor={Louis SCARFONE, Maxence BOURRAGUE}}

% Filet décoratif
\newcommand{\thinrule}{\par\vspace{1mm}{\color{azure!30}\hrule height 0.5pt}\vspace{1mm}}

\begin{document}

% =====================================================================
%  PAGE DE GARDE
% =====================================================================
\begin{titlepage}
\thispagestyle{empty}
\setlength{\parskip}{0pt}
\raggedright

\begin{tcolorbox}[enhanced, sharp corners, boxrule=0pt, colback=azuredark,
   width=\linewidth, left=13mm, right=13mm, top=14mm, bottom=14mm,
   borderline north={5pt}{0pt}{azure}, borderline south={2pt}{0pt}{azure}]
  \raggedright
  {\sffamily\bfseries\large\color{azure!75}\faMicrosoft~~MICROSOFT AZURE \quad\textbullet\quad BLOC 4 — CLOUD COMPUTING}

  \vspace{9mm}
  {\sffamily\bfseries\color{white}\fontsize{37}{42}\selectfont TP2 — Infrastructure as Code}

  \vspace{4mm}
  {\sffamily\color{azurelight}\LARGE avec \textbf{Terraform} pour l'application \textbf{ShopEasy}}

  \vspace{8mm}
  {\color{azure}\rule{5.2cm}{2.5pt}}

  \vspace{5mm}
  {\sffamily\color{white}\large Automatiser le déploiement d'une architecture Azure reproductible, versionnée et maîtrisée par le coût.}
\end{tcolorbox}

\vspace{1.8cm}

\begin{minipage}[t]{0.5\linewidth}
  {\sffamily\bfseries\large\color{azuredark}Binôme}\\[3.5mm]
  {\large Louis \textsc{Scarfone}}\\[2mm]
  {\large Maxence \textsc{Bourragué}}
\end{minipage}\hfill
\begin{minipage}[t]{0.46\linewidth}
  {\sffamily\bfseries\large\color{azuredark}Cadre}\\[3.5mm]
  Mastère Dev Manager Full Stack — RNCP niveau 7\\[2mm]
  Cas fil rouge : \textbf{ShopEasy}\\[2mm]
  Azure for Students — région \texttt{swedencentral}
\end{minipage}

\vspace{1.9cm}

\noindent
\statcard{14}{ateliers + quiz}\hfill
\statcard{21}{ressources Terraform}\hfill
\statcard{$\approx$\,46,5\,\$}{coût mensuel estimé}

\vfill
{\sffamily\large\bfseries\color{azuredark}Juin 2026}\hfill
{\sffamily\color{black!45}Mastère Dev Manager Full Stack}
\end{titlepage}

% =====================================================================
%  TABLE DES MATIÈRES
% =====================================================================
\thispagestyle{fancy}
{\sffamily\Huge\bfseries\color{azuredark}Sommaire}
\vspace{2mm}{\color{azure}\titlerule[1.6pt]}\vspace{4mm}
\hypersetup{linkcolor=azuredark}
{\setlength{\parskip}{1pt}\tableofcontents}
\newpage


% #####################################################################
%  ATELIER 1
% #####################################################################
\section{Atelier 1 — Initialisation du projet Terraform}

\begin{fiche}
\textbf{\eTarget~Objectif} — poser les fondations d'un projet Terraform propre, lisible et sécurisé pour ShopEasy.

\textbf{\ePackage~Livrable attendu} — l'arborescence du projet et le fichier \code{.gitignore}, avec la justification de la non-publication du \code{terraform.tfstate}.
\end{fiche}

Le TP1 a conçu et déployé manuellement (Azure CLI) l'architecture cible de ShopEasy. Le TP2 reprend cette
architecture mais la décrit en \textbf{Infrastructure as Code} avec \textbf{Terraform} : l'infrastructure
devient un ensemble de fichiers \code{.tf} versionnés, relus, appliqués puis détruits de façon reproductible.
Cet atelier met en place le squelette du projet ; aucune ressource Azure n'est encore créée.

\subsection{Vérification de l'environnement}
L'environnement local est vérifié avant toute initialisation :
\begin{codebox}{bash}
\begin{Verbatim}
terraform version
git --version
az version --query '"azure-cli"' -o tsv
az account show --query '{name:name, state:state, user:user.name}' -o json
\end{Verbatim}
\end{codebox}

\begin{codebox}{console}
\begin{Verbatim}
Terraform v1.15.7
on darwin_arm64

git version 2.50.1 (Apple Git-155)

2.87.0

{ "name": "Azure for Students", "state": "Enabled", "user": "louis.scarfone@efrei.net" }
\end{Verbatim}
\end{codebox}

\begin{ltable}{Q[l,0.7] Q[l,0.8] Q[l,0.8] Q[c,0.3]}
Prérequis & Attendu & Constaté & État \\
Terraform CLI & $\geq$ 1.6 & \textbf{v1.15.7} & \eCheck \\
Azure CLI & présent + connecté & \textbf{2.87.0}, session active & \eCheck \\
Git & présent & \textbf{2.50.1} & \eCheck \\
Abonnement Azure & formation, actif & \textbf{Azure for Students}, \textit{Enabled} & \eCheck \\
Clé SSH publique & \code{\textasciitilde/.ssh/id\_rsa.pub} & présente (générée au TP1) & \eCheck \\
\end{ltable}

{\small Terraform n'était pas installé au départ : ajouté via \code{brew install hashicorp/tap/terraform}
— le tap officiel HashiCorp, requis depuis le passage de Terraform sous licence \textbf{BSL}.}

\subsection{Arborescence du projet}
\begin{codebox}{bash}
\begin{Verbatim}
mkdir -p tp2-terraform-azure/templates
cd tp2-terraform-azure
touch versions.tf providers.tf variables.tf locals.tf \
      network.tf security.tf compute.tf loadbalancer.tf \
      storage.tf outputs.tf terraform.tfvars templates/cloud-init.yml
\end{Verbatim}
\end{codebox}

\begin{codebox}{text}
\begin{Verbatim}
tp2-terraform-azure/
|-- .gitignore           exclusions Git (state, secrets, cache)
|-- versions.tf          versions Terraform + providers requis
|-- providers.tf         configuration du provider azurerm
|-- variables.tf         variables d'entree (parametrage du projet)
|-- locals.tf            valeurs calculees (prefixe de nommage, tags communs)
|-- network.tf           Resource Group, VNet, subnet
|-- security.tf          Network Security Group + regles + association
|-- compute.tf           IP publiques, interfaces reseau, 2 VM Linux (cloud-init)
|-- loadbalancer.tf      IP publique LB, Load Balancer, backend pool, sonde, regle
|-- storage.tf           suffixe aleatoire, Storage Account, container Blob
|-- outputs.tf           sorties exposees apres deploiement
|-- terraform.tfvars     valeurs concretes (IGNORE par git -- contient l'IP admin)
`-- templates/
    `-- cloud-init.yml   script d'installation Nginx au premier demarrage des VM
\end{Verbatim}
\end{codebox}

Le projet est découpé par \textbf{responsabilité} (réseau / sécurité / calcul / stockage / sorties) plutôt
qu'en un seul fichier \code{main.tf}. Ce découpage améliore la lisibilité et la maintenabilité : chaque
fichier regroupe un type de ressources cohérent et peut être relu indépendamment. Les fichiers \code{.tf}
sont créés vides à ce stade et seront complétés atelier par atelier.

\subsubsection{Rôle de chaque fichier}
\begin{ltable}{Q[l,0.5] X[l,1]}
Fichier & Rôle \\
\code{versions.tf} & Version minimale de Terraform + providers requis (\code{azurerm}, \code{random}) avec contraintes de version. \\
\code{providers.tf} & Configuration du provider \code{azurerm} (bloc \code{features \{\}}). \\
\code{variables.tf} & Paramètres configurables : projet, environnement, région, admin, clé SSH, CIDR SSH. \\
\code{locals.tf} & Valeurs internes réutilisables : préfixe de nommage, tags communs. \\
\code{network.tf} & Resource Group + Virtual Network + subnet applicatif. \\
\code{security.tf} & NSG, règles entrantes (HTTP, SSH restreint) et association au subnet. \\
\code{compute.tf} & IP publiques, interfaces réseau et 2 VM Linux provisionnées par cloud-init. \\
\code{loadbalancer.tf} & IP publique + Load Balancer L4 (frontend, backend pool, sonde, règle). \\
\code{storage.tf} & Suffixe aléatoire + Storage Account privé + container Blob versionné. \\
\code{outputs.tf} & Sorties utiles (nom du RG, IP du LB, IP des VM, nom du Storage). \\
\code{terraform.tfvars} & Valeurs concrètes de l'environnement \code{dev} (exclu du dépôt). \\
\code{templates/cloud-init.yml} & Configuration cloud-init (installation et démarrage Nginx, page de test). \\
\end{ltable}

\subsection{Fichier \code{.gitignore}}
\begin{codebox}{gitignore}
\begin{Verbatim}
# Repertoire de travail Terraform (providers telecharges, cache, plugins)
.terraform/

# Fichiers de STATE : peuvent contenir des donnees sensibles en clair
*.tfstate
*.tfstate.*

# Variables converties en JSON (peuvent embarquer des valeurs sensibles)
*.tfvars.json

# Valeurs concretes d'environnement : terraform.tfvars contient le CIDR SSH
# = IP publique reelle de l'admin. Seul l'exemple anonymise est versionne.
terraform.tfvars

# Fichier de verrouillage des providers (choix du TP : non versionne)
.terraform.lock.hcl

# Logs de crash Terraform
crash.log
crash.*.log
\end{Verbatim}
\end{codebox}

\begin{ltable}{Q[l,0.45] X[l,1]}
Motif ignoré & Raison \\
\code{.terraform/} & Cache local volumineux (binaires des providers), recréé par \code{terraform init}. \\
\code{*.tfstate} & Le state peut contenir des secrets en clair et l'inventaire complet de l'infrastructure (cf. §4). \\
\code{*.tfvars.json} & Variables au format JSON pouvant contenir des valeurs sensibles. \\
\code{terraform.tfvars} & Contient l'IP publique réelle de l'administrateur — seul l'exemple anonymisé est versionné (§5). \\
\code{.terraform.lock.hcl} & Fichier de verrouillage des providers — non versionné selon le choix du TP (§5). \\
\code{crash.*.log} & Journaux de crash temporaires, sans valeur de suivi. \\
\end{ltable}

\subsection{Point de contrôle — pourquoi ne jamais publier \code{terraform.tfstate}}
Le fichier \code{terraform.tfstate} est une représentation JSON de toutes les ressources gérées par Terraform
et de leurs attributs. Le publier dans un dépôt non sécurisé est dangereux pour quatre raisons :
\begin{enumerate}
  \item \textbf{Présence possible de secrets en clair.} Terraform stocke les attributs renvoyés par Azure, y compris des valeurs sensibles : clés d'accès Storage, mots de passe admin, chaînes de connexion, tokens SAS. Même si aucun secret n'est dans les \code{.tf}, il peut se retrouver dans le state.
  \item \textbf{Exposition de l'inventaire complet} de l'infrastructure (noms, identifiants, IP, plages de subnets, règles réseau) — une cartographie de reconnaissance pour un attaquant.
  \item \textbf{Rôle de source de vérité.} S'il est exposé, corrompu ou modifié, Terraform peut détruire ou recréer des ressources par erreur, ou perdre la trace de ce qu'il gère.
  \item \textbf{Incompatibilité avec le travail d'équipe en local} : conflits de fusion et absence de verrouillage \arr{} corruption du fichier.
\end{enumerate}
La solution professionnelle consiste à externaliser le state dans un \textbf{backend distant} (Storage
Account Azure dédié) avec chiffrement, RBAC et verrouillage — objet de l'Atelier 12. En attendant, le
\code{.gitignore} garantit que le state ne quitte jamais le poste local.

\subsection{Protection des valeurs sensibles}
\textbf{Exclusion de \code{terraform.tfvars}.} Le \code{.gitignore} fourni par le TP n'exclut pas
\code{terraform.tfvars}. Or ce fichier contiendra \code{allowed\_ssh\_cidr} = l'IP publique réelle de
l'administrateur. Conformément à la règle du cours (« ne jamais mettre de valeur sensible dans un fichier
versionné ») et à la pratique du TP1 (masquage de l'IP), \code{terraform.tfvars} est exclu du dépôt ; un
\code{terraform.tfvars.example} anonymisé est versionné à sa place (Atelier 3).

\textbf{Fichier \code{.terraform.lock.hcl}.} Le TP demande de l'ignorer. En projet d'équipe réel, HashiCorp
recommande plutôt de le versionner pour figer les versions de providers et garantir des déploiements
identiques. Le choix du TP est suivi ici, la bonne pratique étant notée.

\subsection{Capture}
\screenshot{tp2-arborescence.png}{Arborescence du projet \code{tp2-terraform-azure} dans l'éditeur : 10 fichiers \code{.tf}, \code{templates/cloud-init.yml} et \code{.gitignore}.}

\begin{etatbox}
\begin{itemize}
  \item Arborescence \code{tp2-terraform-azure/} créée : \textbf{10 fichiers \code{.tf}} (vides), \code{terraform.tfvars}, \code{templates/cloud-init.yml} et \code{.gitignore}.
  \item Environnement validé : Terraform v1.15.7, Azure CLI 2.87.0 (\textit{Azure for Students}), Git 2.50.1.
  \item \code{.gitignore} en place : state, secrets et \code{terraform.tfvars} (IP admin) exclus du dépôt ; point de contrôle \code{tfstate} traité.
  \item \textit{Région \code{francecentral} et taille \code{B1s} du TP non disponibles sur Azure for Students \arr{} adaptées en \code{swedencentral} et \code{B2ats\_v2} (Ateliers 3 et 6).}
  \item \textbf{Prêt pour l'Atelier 2 — déclaration des providers.}
\end{itemize}
\end{etatbox}


% #####################################################################
%  ATELIER 2
% #####################################################################
\section{Atelier 2 — Déclaration des providers}

\begin{fiche}
\textbf{\eTarget~Objectif} — déclarer la version de Terraform et les providers requis, puis initialiser le projet.

\textbf{\ePackage~Livrable attendu} — \code{versions.tf}, \code{providers.tf} et le résultat de \code{terraform validate}.
\end{fiche}

\subsection{\code{versions.tf}}
\begin{codebox}{hcl}
\begin{Verbatim}
terraform {
  required_version = ">= 1.6.0"

  required_providers {
    azurerm = {
      source  = "hashicorp/azurerm"
      version = "~> 4.0"
    }
    random = {
      source  = "hashicorp/random"
      version = "~> 3.6"
    }
  }
}
\end{Verbatim}
\end{codebox}

\begin{ltable}{Q[l,0.5] X[l,1]}
Élément & Signification \\
\code{required\_version = ">= 1.6.0"} & Le projet exige Terraform 1.6 ou plus (version installée : 1.15.7). \\
\code{azurerm} & Provider standard pour Azure : RG, VNet, NSG, VM, Load Balancer, Storage. \\
\code{random} & Génère un suffixe aléatoire pour rendre le nom du Storage Account globalement unique (Atelier 8). \\
\code{version = "\textasciitilde> 4.0"} & Opérateur pessimiste : autorise \code{>= 4.0.0} et \code{< 5.0.0} (toutes les \code{4.x}). \\
\code{version = "\textasciitilde> 3.6"} & Autorise \code{>= 3.6.0} et \code{< 4.0.0}. \\
\end{ltable}

L'opérateur \code{\textasciitilde>} fige la branche majeure d'un provider : les correctifs \code{4.x} sont
récupérés automatiquement, mais une version \code{5.0} (susceptible d'introduire des \textit{breaking
changes}) est exclue.

\subsection{\code{providers.tf}}
\begin{codebox}{hcl}
\begin{Verbatim}
provider "azurerm" {
  features {}
}
\end{Verbatim}
\end{codebox}

Le bloc \code{features \{\}} est obligatoire pour le provider \code{azurerm}, même vide : il active le
comportement par défaut du provider. L'authentification est héritée de la session Azure CLI (\code{az login})
déjà active, ce qui évite d'inscrire un identifiant dans le code.

\begin{tipbox}[Authentification azurerm v4.x]
Le provider \code{azurerm} v4.x requiert un \code{subscription\_id} lors des opérations \code{plan} et
\code{apply} (pas pour \code{init} ni \code{validate}). Il sera fourni via la variable d'environnement
\code{ARM\_SUBSCRIPTION\_ID} à l'Atelier 4, pour conserver \code{providers.tf} minimal et sans identifiant en
dur.
\end{tipbox}

\subsection{Initialisation — \code{terraform init}}
\begin{codebox}{console}
\begin{Verbatim}
Initializing provider plugins found in the configuration...
- Finding hashicorp/random versions matching "~> 3.6"...
- Finding hashicorp/azurerm versions matching "~> 4.0"...
- Installing hashicorp/random v3.9.0...
- Installed hashicorp/random v3.9.0 (signed by HashiCorp)
- Installing hashicorp/azurerm v4.78.0...
- Installed hashicorp/azurerm v4.78.0 (signed by HashiCorp)

Terraform has created a lock file .terraform.lock.hcl to record the provider
selections it made above.

Terraform has been successfully initialized!
\end{Verbatim}
\end{codebox}
\code{terraform init} a résolu et téléchargé les providers (vérifiés par signature HashiCorp) :
\code{azurerm v4.78.0} et \code{random v3.9.0}. La commande a créé le dossier de travail \code{.terraform/}
et le fichier \code{.terraform.lock.hcl} qui fige les versions exactes retenues.

\subsection{Formatage — \code{terraform fmt}}
La commande ne renvoie aucune sortie : tous les fichiers \code{.tf} sont déjà conformes au style canonique
HCL. \code{terraform fmt} n'affiche que les fichiers qu'il modifie ; une sortie vide est le résultat attendu.

\subsection{Validation — \code{terraform validate}}
\begin{codebox}{console}
\begin{Verbatim}
Success! The configuration is valid.
\end{Verbatim}
\end{codebox}
La configuration est syntaxiquement et structurellement valide. \code{terraform validate} travaille en local :
il vérifie la cohérence du code sans contacter Azure ni l'existence réelle des ressources.

\begin{etatbox}
\begin{itemize}
  \item \code{versions.tf} : Terraform \code{>= 1.6.0}, providers \textbf{azurerm \code{\textasciitilde> 4.0}} et \textbf{random \code{\textasciitilde> 3.6}}.
  \item \code{providers.tf} : provider \code{azurerm} configuré (\code{features \{\}}), authentification héritée d'\code{az login}.
  \item \code{terraform init} réussi : \code{azurerm v4.78.0} et \code{random v3.9.0} installés, lock file créé ; \code{validate} : \textbf{Success}.
  \item \textbf{Prêt pour l'Atelier 3 — paramétrage du projet.}
\end{itemize}
\end{etatbox}

% #####################################################################
%  ATELIER 3
% #####################################################################
\section{Atelier 3 — Paramétrage du projet}

\begin{fiche}
\textbf{\eTarget~Objectif} — rendre le projet réutilisable en externalisant les valeurs configurables et en centralisant le nommage et les tags.

\textbf{\ePackage~Livrable attendu} — \code{variables.tf}, \code{terraform.tfvars} et \code{locals.tf} (variables, valeurs d'environnement, locals et tags).
\end{fiche}

\subsection{Variables d'entrée — \code{variables.tf}}
\begin{codebox}{hcl}
\begin{Verbatim}
variable "project" {
  description = "Nom court du projet"
  type        = string
  default     = "shopeasy"
}

variable "environment" {
  description = "Nom de l'environnement"
  type        = string
  default     = "dev"
}

variable "location" {
  description = "Region Azure cible"
  type        = string
  default     = "francecentral"
}

variable "admin_username" {
  description = "Utilisateur administrateur Linux"
  type        = string
  default     = "azureuser"
}

variable "ssh_public_key_path" {
  description = "Chemin local vers la cle publique SSH"
  type        = string
}

variable "allowed_ssh_cidr" {
  description = "CIDR autorise pour l'acces SSH"
  type        = string
}
\end{Verbatim}
\end{codebox}

\begin{ltable}{Q[l,0.55] Q[c,0.35] Q[l,0.45] X[l,1]}
Variable & Type & Défaut & Rôle \\
\code{project} & string & \code{shopeasy} & Nom court du projet ; alimente le préfixe et les tags. \\
\code{environment} & string & \code{dev} & Environnement ; alimente le préfixe et les tags. \\
\code{location} & string & \code{francecentral} & Région Azure cible (surchargée par \code{terraform.tfvars}, §2). \\
\code{admin\_username} & string & \code{azureuser} & Compte administrateur Linux des VM. \\
\code{ssh\_public\_key\_path} & string & \textit{(requis)} & Chemin de la clé publique SSH injectée dans les VM. \\
\code{allowed\_ssh\_cidr} & string & \textit{(requis)} & Plage CIDR autorisée en SSH (IP admin en \code{/32}). \\
\end{ltable}

Les variables \code{ssh\_public\_key\_path} et \code{allowed\_ssh\_cidr} sont \textbf{sans valeur par défaut} :
elles sont propres à l'environnement et doivent être fournies via \code{terraform.tfvars}. Externaliser ces
paramètres évite tout codage en dur et permet de réutiliser le même code pour d'autres environnements.

\subsection{Valeurs de l'environnement — \code{terraform.tfvars}}

\begin{notebox}[Adaptation de la région (Azure for Students)]
Le TP fixe \code{location = "francecentral"}. Cette région est \textbf{bloquée} par la policy \textit{Allowed
resource deployment regions} de l'abonnement \textit{Azure for Students} (contrainte déjà rencontrée au TP1).
La région déployable retenue est \textbf{\code{swedencentral}}. La valeur est définie dans
\code{terraform.tfvars} — chargé automatiquement par Terraform — sans modifier le défaut de \code{variables.tf}.
\end{notebox}

Valeurs retenues (modèle anonymisé \code{terraform.tfvars.example}) :
\begin{codebox}{hcl}
\begin{Verbatim}
project             = "shopeasy"
environment         = "dev"
location            = "swedencentral"   # francecentral bloquee sur Azure for Students
admin_username      = "azureuser"
ssh_public_key_path = "~/.ssh/id_rsa.pub"
allowed_ssh_cidr    = "X.X.X.X/32"      # IP publique de l'administrateur (masquee)
\end{Verbatim}
\end{codebox}

Le fichier réel \code{terraform.tfvars} renseigne \code{allowed\_ssh\_cidr} avec l'\textbf{IP publique réelle}
de l'administrateur, en \code{/32}. Conformément à la règle « aucune valeur sensible dans un fichier
versionné », \code{terraform.tfvars} est \textbf{exclu du dépôt} ; seul \code{terraform.tfvars.example},
anonymisé, est versionné.

\subsection{Locals et tags — \code{locals.tf}}
\begin{codebox}{hcl}
\begin{Verbatim}
locals {
  # Prefixe de nommage commun a toutes les ressources : "shopeasy-dev"
  prefix = "${var.project}-${var.environment}"

  # Tags de gouvernance appliques a chaque ressource
  common_tags = {
    project     = var.project
    environment = var.environment
    owner       = "formation"
    managed_by  = "terraform"
  }
}
\end{Verbatim}
\end{codebox}

\code{prefix} vaut \code{shopeasy-dev} et est réutilisé pour nommer chaque ressource de façon cohérente
(\code{rg-shopeasy-dev}, \code{vnet-shopeasy-dev}, \code{nsg-shopeasy-dev-web}…). \code{common\_tags} regroupe
les tags de gouvernance appliqués à toutes les ressources :

\begin{ltable}{Q[l,0.4] Q[l,0.5] X[l,1]}
Tag & Valeur & Usage \\
\code{project} & \code{shopeasy} & Rattachement applicatif. \\
\code{environment} & \code{dev} & Distinction dev / test / prod. \\
\code{owner} & \code{formation} & Équipe responsable. \\
\code{managed\_by} & \code{terraform} & Indique que la ressource est gérée en IaC (toute modification manuelle est une dérive). \\
\end{ltable}

Centraliser le nommage et les tags dans \code{locals.tf} applique le principe \textit{DRY} : une seule source
de vérité, modifiable en un point, propagée à l'ensemble des ressources.

\subsection{Validation}
\begin{codebox}{console}
\begin{Verbatim}
$ terraform fmt        # aucune sortie : fichiers deja conformes
$ terraform validate
Success! The configuration is valid.

# Controle de la protection de l'IP par .gitignore
terraform.tfvars          -> IGNORE par git (OK, contient l'IP)
terraform.tfvars.example  -> VERSIONNE (OK, anonymise)
\end{Verbatim}
\end{codebox}
La configuration est valide et le fichier contenant l'IP réelle est bien exclu du dépôt.

\begin{etatbox}
\begin{itemize}
  \item \code{variables.tf} : 6 variables d'entrée (projet, environnement, région, admin, clé SSH, CIDR SSH).
  \item \code{terraform.tfvars} : environnement \code{dev}, \textbf{région \code{swedencentral}} (adaptation Azure for Students), IP admin en \code{/32} — fichier \textbf{exclu du dépôt}.
  \item \code{terraform.tfvars.example} : modèle \textbf{anonymisé} versionné ; \code{locals.tf} : préfixe \code{shopeasy-dev} + 4 tags de gouvernance.
  \item \textbf{Prêt pour l'Atelier 4 — Resource Group et réseau (premier \code{apply}).}
\end{itemize}
\end{etatbox}


% #####################################################################
%  ATELIER 4
% #####################################################################
\section{Atelier 4 — Groupe de ressources et réseau}

\begin{fiche}
\textbf{\eTarget~Objectif} — créer le groupe de ressources et le réseau privé de ShopEasy avec Terraform, et réaliser le premier \code{apply}.

\textbf{\ePackage~Livrable attendu} — \code{network.tf} (RG + VNet + subnet) et la preuve de création (plan, apply, portail).

\textbf{\eGlobe~Région} — \code{swedencentral} \quad\textbullet\quad RG \code{rg-shopeasy-dev}, VNet \code{10.20.0.0/16}.
\end{fiche}

\subsection{Authentification du provider}
L'authentification vers Azure est héritée de la session Azure CLI (\code{az login}). Le provider
\code{azurerm} v4.x exige en plus l'identifiant d'abonnement, fourni par variable d'environnement, \textbf{sans
l'inscrire dans le code} :
\begin{codebox}{bash}
\begin{Verbatim}
export ARM_SUBSCRIPTION_ID=$(az account show --query id -o tsv)
\end{Verbatim}
\end{codebox}

\subsection{Ressources réseau — \code{network.tf}}
\begin{codebox}{hcl}
\begin{Verbatim}
resource "azurerm_resource_group" "main" {
  name     = "rg-${local.prefix}"
  location = var.location
  tags     = local.common_tags
}

resource "azurerm_virtual_network" "main" {
  name                = "vnet-${local.prefix}"
  address_space       = ["10.20.0.0/16"]
  location            = azurerm_resource_group.main.location
  resource_group_name = azurerm_resource_group.main.name
  tags                = local.common_tags
}

resource "azurerm_subnet" "web" {
  name                 = "snet-web"
  resource_group_name  = azurerm_resource_group.main.name
  virtual_network_name = azurerm_virtual_network.main.name
  address_prefixes     = ["10.20.1.0/24"]
}
\end{Verbatim}
\end{codebox}

Le VNet et le subnet \textbf{référencent} les attributs du Resource Group : Terraform en déduit le
\textbf{graphe de dépendances} (RG d'abord, puis VNet, puis subnet). Le nommage s'appuie sur \code{local.prefix}
et chaque élément porte \code{local.common\_tags}.

\begin{ltable}{Q[l,0.9] Q[l,0.7] X[l,0.9]}
Ressource & Nom & Valeur clé \\
\code{azurerm\_resource\_group.main} & \code{rg-shopeasy-dev} & région \code{swedencentral} \\
\code{azurerm\_virtual\_network.main} & \code{vnet-shopeasy-dev} & espace \code{10.20.0.0/16} \\
\code{azurerm\_subnet.web} & \code{snet-web} & préfixe \code{10.20.1.0/24} \\
\end{ltable}

\subsection{Prévisualisation — \code{terraform plan}}
\begin{codebox}{console}
\begin{Verbatim}
Terraform will perform the following actions:

  # azurerm_resource_group.main will be created
  + resource "azurerm_resource_group" "main" {
      + id       = (known after apply)
      + location = "swedencentral"
      + name     = "rg-shopeasy-dev"
      + tags     = {
          + "environment" = "dev"
          + "managed_by"  = "terraform"
          + "owner"       = "formation"
          + "project"     = "shopeasy"
        }
    }

  # azurerm_subnet.web will be created
  + resource "azurerm_subnet" "web" {
      + address_prefixes    = [ "10.20.1.0/24" ]
      + name                = "snet-web"
      + resource_group_name = "rg-shopeasy-dev"
      + virtual_network_name = "vnet-shopeasy-dev"
    }

  # azurerm_virtual_network.main will be created
  + resource "azurerm_virtual_network" "main" {
      + address_space       = [ "10.20.0.0/16" ]
      + location            = "swedencentral"
      + name                = "vnet-shopeasy-dev"
      + resource_group_name = "rg-shopeasy-dev"
      + tags                = { ... 4 tags ... }
    }

Plan: 3 to add, 0 to change, 0 to destroy.
\end{Verbatim}
\end{codebox}
Le plan annonce \textbf{3 ressources à créer}, aucune modification ni destruction — résultat attendu pour un
premier déploiement.

\subsection{Application — \code{terraform apply}}
\begin{codebox}{console}
\begin{Verbatim}
azurerm_resource_group.main: Creating...
azurerm_resource_group.main: Creation complete after 26s [id=/subscriptions/e67b15f9-.../resourceGroups/rg-shopeasy-dev]
azurerm_virtual_network.main: Creating...
azurerm_virtual_network.main: Creation complete after 6s [id=.../virtualNetworks/vnet-shopeasy-dev]
azurerm_subnet.web: Creating...
azurerm_subnet.web: Creation complete after 5s [id=.../subnets/snet-web]

Apply complete! Resources: 3 added, 0 changed, 0 destroyed.
\end{Verbatim}
\end{codebox}
{\small L'identifiant d'abonnement est masqué (\code{e67b15f9-...}) dans les ID de ressources.}

\subsection{Vérification (Azure CLI et état Terraform)}
\begin{codebox}{console}
\begin{Verbatim}
$ az group show -n rg-shopeasy-dev --query '{name,location,provisioningState:...,tags}'
{ "location": "swedencentral", "name": "rg-shopeasy-dev",
  "provisioningState": "Succeeded",
  "tags": { "environment":"dev","managed_by":"terraform","owner":"formation","project":"shopeasy" } }

$ az network vnet show -g rg-shopeasy-dev -n vnet-shopeasy-dev --query '{...}'
{ "addressSpace": [ "10.20.0.0/16" ], "location": "swedencentral",
  "name": "vnet-shopeasy-dev", "provisioningState": "Succeeded" }
  -> subnet snet-web : addressPrefix 10.20.1.0/24, state Succeeded

$ terraform state list
azurerm_resource_group.main
azurerm_subnet.web
azurerm_virtual_network.main
\end{Verbatim}
\end{codebox}
Les trois ressources sont en \code{provisioningState: Succeeded}, avec la bonne région, le bon plan
d'adressage et les quatre tags de gouvernance. L'état Terraform recense exactement les trois ressources gérées.

\subsection{Captures portail}
\screenshot{tp2-rg-overview.png}{Vue d'ensemble du Resource Group \code{rg-shopeasy-dev} (région Sweden Central) — \textit{subscription ID flouté}.}
\screenshot{tp2-rg-tags.png}{Les 4 tags de gouvernance appliqués au Resource Group.}
\screenshot{tp2-vnet-subnets.png}{Le subnet \code{snet-web} (\code{10.20.1.0/24}) dans le VNet \code{vnet-shopeasy-dev}.}

\subsection{Questions}
\textbf{1. Pourquoi le VNet utilise-t-il une plage privée ?}\\
La plage \code{10.20.0.0/16} appartient aux \textbf{adresses privées RFC 1918}, non routables sur Internet. Un
VNet est un \textbf{réseau privé} : une plage privée évite tout conflit avec des adresses publiques, garde les
communications internes \textbf{isolées}, et permet une \textbf{segmentation maîtrisée} par subnets.
L'exposition vers l'extérieur se fait \textbf{explicitement} via des IP publiques et un Load Balancer
(Ateliers 6 et 7), jamais par le VNet lui-même.

\textbf{2. Que faudrait-il ajouter pour isoler une base de données dans un subnet dédié ?}\\
Un \textbf{second subnet} (ex. \code{snet-data}, \code{10.20.2.0/24}) réservé à la couche données, associé à
son \textbf{propre NSG} n'autorisant l'accès à la base (port \code{1433}/\code{5432}) \textbf{que depuis le
subnet web} et refusant le reste. Idéalement, la base est jointe via un \textbf{Private Endpoint} (aucune
exposition publique). C'est l'application de la \textbf{segmentation réseau} et du \textbf{moindre privilège}.

\textbf{3. Pourquoi séparer les fichiers Terraform au lieu de tout placer dans \code{main.tf} ?}\\
Terraform charge \textbf{tous} les fichiers \code{.tf} d'un répertoire quel que soit leur nom : le découpage
est purement \textbf{organisationnel}, sans coût fonctionnel. Regrouper les ressources par \textbf{responsabilité}
rend le projet \textbf{lisible}, \textbf{navigable} et \textbf{maintenable}, et réduit les conflits de fusion
en équipe.

\begin{etatbox}
\begin{itemize}
  \item \code{network.tf} créé ; \code{terraform apply} : \textbf{3 ressources ajoutées} (\code{rg-shopeasy-dev}, \code{vnet-shopeasy-dev}, \code{snet-web}), toutes en \code{Succeeded}.
  \item Région \code{swedencentral}, plan d'adressage \code{10.20.0.0/16} (subnet web \code{10.20.1.0/24}), 4 tags de gouvernance.
  \item \textbf{Prêt pour l'Atelier 5 — sécurisation du réseau avec un NSG.}
\end{itemize}
\end{etatbox}

% #####################################################################
%  ATELIER 5
% #####################################################################
\section{Atelier 5 — Sécurisation du réseau avec un NSG}

\begin{fiche}
\textbf{\eTarget~Objectif} — filtrer les flux réseau du subnet web en n'autorisant que les ouvertures nécessaires.

\textbf{\ePackage~Livrable attendu} — \code{security.tf} (NSG + règles + association) et l'analyse de sécurité des flux.
\end{fiche}

\subsection{NSG, règles et association — \code{security.tf}}
\begin{codebox}{hcl}
\begin{Verbatim}
resource "azurerm_network_security_group" "web" {
  name                = "nsg-${local.prefix}-web"
  location            = azurerm_resource_group.main.location
  resource_group_name = azurerm_resource_group.main.name
  tags                = local.common_tags

  security_rule {
    name                       = "Allow-HTTP"
    priority                   = 100
    direction                  = "Inbound"
    access                     = "Allow"
    protocol                   = "Tcp"
    source_port_range          = "*"
    destination_port_range     = "80"
    source_address_prefix      = "Internet"
    destination_address_prefix = "*"
  }

  security_rule {
    name                       = "Allow-SSH-Admin"
    priority                   = 110
    direction                  = "Inbound"
    access                     = "Allow"
    protocol                   = "Tcp"
    source_port_range          = "*"
    destination_port_range     = "22"
    source_address_prefix      = var.allowed_ssh_cidr
    destination_address_prefix = "*"
  }
}

resource "azurerm_subnet_network_security_group_association" "web" {
  subnet_id                 = azurerm_subnet.web.id
  network_security_group_id = azurerm_network_security_group.web.id
}
\end{Verbatim}
\end{codebox}

\begin{ltable}{Q[l,0.55] Q[c,0.3] Q[c,0.3] Q[c,0.25] Q[c,0.25] Q[l,0.55] X[l,0.9]}
Règle & Prio & Sens & Accès & Port & Source & Rôle \\
\code{Allow-HTTP} & 100 & In & Allow & 80 & \code{Internet} & Accès HTTP à l'application web de test \\
\code{Allow-SSH-Admin} & 110 & In & Allow & 22 & \code{<IP\_ADMIN>/32} & Administration Linux limitée à l'IP admin \\
\end{ltable}

La source SSH n'est jamais codée en dur : elle provient de la variable \code{allowed\_ssh\_cidr} (définie dans
\code{terraform.tfvars}, exclu du dépôt). L'association rattache le NSG au subnet \code{snet-web} : toutes les
VM qui y seront déployées héritent du filtrage.

\subsection{Prévisualisation et application}
Le plan annonce \textbf{2 ressources à créer} (NSG + association). Extrait (IP admin masquée) :
\begin{codebox}{console}
\begin{Verbatim}
  # azurerm_network_security_group.web will be created
  + security_rule = [
      + { name = "Allow-SSH-Admin", priority = 110, direction = "Inbound",
          access = "Allow", protocol = "Tcp", destination_port_range = "22",
          source_address_prefix = "<IP_ADMIN>/32" },
      + { name = "Allow-HTTP", priority = 100, direction = "Inbound",
          access = "Allow", protocol = "Tcp", destination_port_range = "80",
          source_address_prefix = "Internet" },
    ]

Plan: 2 to add, 0 to change, 0 to destroy.
azurerm_network_security_group.web: Creation complete after 2s [id=.../nsg-shopeasy-dev-web]
azurerm_subnet_network_security_group_association.web: Creation complete after 6s
Apply complete! Resources: 2 added, 0 changed, 0 destroyed.
\end{Verbatim}
\end{codebox}

\subsection{Vérification (Azure CLI)}
\begin{codebox}{console}
\begin{Verbatim}
$ az network nsg rule list -g rg-shopeasy-dev --nsg-name nsg-shopeasy-dev-web ...
Name             Priority  Direction  Access  Proto  Port  Source
Allow-HTTP       100       Inbound    Allow   Tcp    80    Internet
Allow-SSH-Admin  110       Inbound    Allow   Tcp    22    <IP_ADMIN>/32

$ az network vnet subnet show ... -n snet-web --query '{subnet,nsg_associe}'
{ "nsg_associe": ".../networkSecurityGroups/nsg-shopeasy-dev-web", "subnet": "snet-web" }
\end{Verbatim}
\end{codebox}
Les deux règles sont actives et le subnet \code{snet-web} est associé au NSG \code{nsg-shopeasy-dev-web}. La
source de la règle SSH est restreinte à l'IP de l'administrateur (masquée).

\subsection{Captures portail}
\screenshot{tp2-nsg-rules.png}{Règles entrantes du NSG \code{nsg-shopeasy-dev-web} (\code{Allow-HTTP} et \code{Allow-SSH-Admin}) — \textit{source SSH (IP admin) floutée}.}
\screenshot{tp2-nsg-association.png}{Association du NSG \code{nsg-shopeasy-dev-web} au subnet \code{snet-web}.}

\subsection{Analyse de sécurité des flux}
\begin{tblr}{
  width=\linewidth, colspec={Q[l,0.6] Q[c,0.3] X[l,1.1] X[l,1.1]},
  row{1}={bg=azuredark,fg=white,font=\sffamily\bfseries\scriptsize},
  row{2-Z}={font=\scriptsize}, row{even}={bg=rowalt},
  hline{1,Z}={1.1pt,azuredark}, hline{2}={0.5pt,azuredark!45},
  colsep=4pt, rowsep=3.2pt, stretch=1.15,
}
Flux & Autorisé ? & Justification & Risque résiduel \\
\textbf{Internet \arr{} HTTP (80)} & \textbf{Oui} & Règle \code{Allow-HTTP} ; expose l'application web de test. & Trafic \textbf{en clair}. En prod : HTTPS (443) + Application Gateway + WAF. \\
\textbf{SSH (22) depuis l'IP admin} & \textbf{Oui} & Règle \code{Allow-SSH-Admin} limitée à l'IP admin en \code{/32}. & Dépend d'une IP fixe. En prod : Azure Bastion / VPN, JIT, MFA. \\
\textbf{SSH (22) depuis tout Internet} & \textbf{Non} & Aucune règle ne l'autorise ; \code{DenyAllInBound} (65500) bloque le reste. & Quasi nul tant que le CIDR admin reste étroit (\code{/32}). \\
\textbf{Tout trafic sortant} & \textbf{Oui} & Azure autorise l'Outbound par défaut (\code{AllowVnetOutBound}, \code{AllowInternetOutBound}). & Exfiltration possible. En prod : restreindre l'Outbound aux destinations légitimes. \\
\end{tblr}
{\small Un NSG applique des règles par défaut invisibles : \code{AllowVnetInBound}, \code{AllowAzureLoadBalancerInBound}, \code{DenyAllInBound} en entrée. Les règles personnalisées (prio < 65000) sont évaluées \textbf{avant} ces défauts.}

\subsection{Question de recul — l'accès SSH en production}
Exposer le port SSH (22) directement sur Internet, même restreint à une IP, reste à \textbf{éviter en
production}. Les approches recommandées :
\begin{itemize}
  \item \textbf{Azure Bastion} : connexion SSH/RDP via le portail, \textbf{sans IP publique} sur les VM ;
  \item \textbf{VPN} site-to-site ou point-to-site : administration par le réseau privé ;
  \item \textbf{Jump server / bastion host} contrôlé et audité ;
  \item \textbf{Accès just-in-time} (Defender for Cloud) : ouverture à la demande, pour une durée et une IP limitées.
\end{itemize}
L'objectif commun est de \textbf{supprimer l'exposition permanente} du SSH au profit d'un accès privé,
temporaire et tracé.

\begin{etatbox}
\begin{itemize}
  \item \code{security.tf} créé : NSG \code{nsg-shopeasy-dev-web} (2 règles entrantes) + association au subnet \code{snet-web} ; \code{apply} : \textbf{2 ressources ajoutées}.
  \item HTTP (80) ouvert depuis Internet ; SSH (22) \textbf{restreint à l'IP admin} ; reste de l'entrant bloqué par défaut.
  \item \textbf{Prêt pour l'Atelier 6 — déploiement des deux machines virtuelles Linux.}
\end{itemize}
\end{etatbox}


% #####################################################################
%  ATELIER 6
% #####################################################################
\section{Atelier 6 — Déploiement des deux machines virtuelles Linux}

\begin{fiche}
\textbf{\eTarget~Objectif} — déployer deux serveurs web Linux identiques, provisionnés automatiquement par cloud-init.

\textbf{\ePackage~Livrable attendu} — \code{templates/cloud-init.yml}, \code{compute.tf} (IP publiques, interfaces réseau, VM) et la preuve d'accès web.
\end{fiche}

\begin{notebox}[Adaptation imposée par Azure for Students (taille de VM)]
Le TP fixe \code{size = "Standard\_B1s"}, \textbf{indisponible} sur l'abonnement \textit{Azure for Students}
(contrainte déjà résolue au TP1). La taille retenue est \textbf{\code{Standard\_B2ats\_v2}} (2 vCPU, série
burstable \code{Basv2}), \textbf{officiellement recommandée par Microsoft} comme alternative à \code{B1s}. Le
quota \code{Basv2} (10 vCPU) couvre les 2 × 2 = 4 vCPU nécessaires.
\end{notebox}

\subsection{Script cloud-init — \code{templates/cloud-init.yml}}
\begin{codebox}{yaml}
\begin{Verbatim}
#cloud-config
package_update: true
packages:
  - nginx
write_files:
  - path: /var/www/html/index.html
    content: |
      <html>
        <body>
          <h1>ShopEasy - serveur web ${server_index}</h1>
          <p>Deploiement realise avec Terraform sur Azure.</p>
        </body>
      </html>
runcmd:
  - systemctl enable nginx
  - systemctl restart nginx
\end{Verbatim}
\end{codebox}
Ce script est exécuté \textbf{au premier démarrage} de chaque VM : il installe \textbf{Nginx}, écrit une page
\code{index.html} puis active et démarre le service. Le marqueur \code{\$\{server\_index\}} est injecté par
Terraform (§2) pour \textbf{personnaliser la page par serveur}.

\subsection{IP publiques, interfaces et VM — \code{compute.tf}}
\begin{codebox}{hcl}
\begin{Verbatim}
resource "azurerm_public_ip" "web" {
  count               = 2
  name                = "pip-${local.prefix}-web-${count.index + 1}"
  location            = azurerm_resource_group.main.location
  resource_group_name = azurerm_resource_group.main.name
  allocation_method   = "Static"
  sku                 = "Standard"
  tags                = local.common_tags
}

resource "azurerm_network_interface" "web" {
  count               = 2
  name                = "nic-${local.prefix}-web-${count.index + 1}"
  location            = azurerm_resource_group.main.location
  resource_group_name = azurerm_resource_group.main.name
  tags                = local.common_tags

  ip_configuration {
    name                          = "internal"
    subnet_id                     = azurerm_subnet.web.id
    private_ip_address_allocation = "Dynamic"
    public_ip_address_id          = azurerm_public_ip.web[count.index].id
  }
}

resource "azurerm_linux_virtual_machine" "web" {
  count               = 2
  name                = "vm-${local.prefix}-web-${count.index + 1}"
  location            = azurerm_resource_group.main.location
  resource_group_name = azurerm_resource_group.main.name
  size                = "Standard_B2ats_v2"   # B1s indisponible (Azure for Students)
  admin_username      = var.admin_username
  network_interface_ids = [ azurerm_network_interface.web[count.index].id ]

  admin_ssh_key {
    username   = var.admin_username
    public_key = file(pathexpand(var.ssh_public_key_path))
  }

  os_disk {
    caching              = "ReadWrite"
    storage_account_type = "Standard_LRS"
  }

  source_image_reference {
    publisher = "Canonical"
    offer     = "0001-com-ubuntu-server-jammy"
    sku       = "22_04-lts"
    version   = "latest"
  }

  custom_data = base64encode(templatefile("${path.module}/templates/cloud-init.yml", {
    server_index = count.index + 1
  }))

  tags = local.common_tags
}
\end{Verbatim}
\end{codebox}

\begin{itemize}
  \item \textbf{\code{count = 2}} déploie deux exemplaires de chaque ressource, indexés par \code{count.index} (0 et 1). Chaque VM est reliée à \textbf{sa} NIC et \textbf{son} IP via \code{[count.index]}.
  \item \textbf{\code{templatefile(...)}} lit le \code{cloud-init.yml} et y injecte \code{server\_index = count.index + 1}, puis \code{base64encode} encode le résultat pour \code{custom\_data}.
  \item \textbf{\code{pathexpand(...)}} : la fonction \code{file()} ne développe pas \code{\textasciitilde} ; \code{pathexpand} convertit \code{\textasciitilde/.ssh/id\_rsa.pub} en chemin absolu.
\end{itemize}

\subsection{Plan et application}
\begin{codebox}{console}
\begin{Verbatim}
Plan: 6 to add, 0 to change, 0 to destroy.   (2 IP publiques, 2 NIC, 2 VM)

azurerm_public_ip.web[0..1]: Creation complete after ~3s
azurerm_network_interface.web[0..1]: Creation complete after ~5s
azurerm_linux_virtual_machine.web[1]: Creation complete after 49s [id=.../vm-shopeasy-dev-web-2]
azurerm_linux_virtual_machine.web[0]: Creation complete after 50s [id=.../vm-shopeasy-dev-web-1]

Apply complete! Resources: 6 added, 0 changed, 0 destroyed.
\end{Verbatim}
\end{codebox}

\subsection{Vérification (Azure CLI et accès web)}
\begin{ltable}{Q[l,0.9] Q[l,0.8] Q[l,0.6] Q[l,0.5] Q[c,0.5]}
Nom & Taille & IP publique & IP privée & État \\
\code{vm-shopeasy-dev-web-1} & \code{Standard\_B2ats\_v2} & \code{135.225.73.252} & \code{10.20.1.5} & VM running \\
\code{vm-shopeasy-dev-web-2} & \code{Standard\_B2ats\_v2} & \code{135.225.74.4} & \code{10.20.1.4} & VM running \\
\end{ltable}

Test HTTP de chaque page (depuis le poste admin) :
\begin{codebox}{console}
\begin{Verbatim}
$ curl http://135.225.73.252
<html><body><h1>ShopEasy - serveur web 1</h1>
<p>Deploiement realise avec Terraform sur Azure.</p></body></html>

$ curl http://135.225.74.4
<html><body><h1>ShopEasy - serveur web 2</h1>
<p>Deploiement realise avec Terraform sur Azure.</p></body></html>
\end{Verbatim}
\end{codebox}
Les deux VM sont \code{VM running}, en \code{Standard\_B2ats\_v2}, et chaque page identifie son serveur (1 et
2) — ce qui permettra de vérifier la répartition de charge à l'Atelier 7. L'état Terraform recense les
\textbf{11 ressources} gérées.

\subsection{Captures portail et navigateur}
\screenshot{tp2-vms-in-rg.png}{Les deux VM \code{vm-shopeasy-dev-web-1/2} (\code{VM running}, \code{Standard\_B2ats\_v2}).}
\screenshot{tp2-web-vm1.png}{Page web servie par \code{vm-web-1} (\code{http://135.225.73.252}) : « serveur web 1 ».}
\screenshot{tp2-web-vm2.png}{Page web servie par \code{vm-web-2} (\code{http://135.225.74.4}) : « serveur web 2 ».}

\subsection{Questions}
\textbf{1. À quoi sert \code{count} dans cette configuration ?}\\
\code{count = 2} crée \textbf{deux exemplaires} de la ressource, indexés par \code{count.index} (0 et 1). Cela
permet de déployer plusieurs serveurs \textbf{sans dupliquer le code} (DRY). L'index sert à \textbf{nommer}
les ressources (\code{web-1}, \code{web-2}) et à \textbf{relier} chaque VM à la bonne interface et à la bonne IP.

\textbf{2. Quel est le rôle de \code{custom\_data} ?}\\
\code{custom\_data} transmet à la VM un script \textbf{cloud-init} (encodé en base64) exécuté \textbf{au premier
démarrage}. Ici, il installe Nginx et publie la page de test, sans aucune intervention manuelle ni connexion
SSH. Couplé à \code{templatefile}, il \textbf{personnalise} la page de chaque serveur.

\textbf{3. Pourquoi une VM \code{Standard\_B2ats\_v2} (à défaut de \code{B1s}) est-elle acceptable en
formation ?}\\
Une VM \textbf{burstable de petite taille} (série B) suffit à un environnement de \textbf{test/formation} :
charge faible, pas de SLA de production, \textbf{coût minimal} ($\approx$ 0,01 \$/h). Sur \textit{Azure for
Students}, \code{B1s} est indisponible : \code{B2ats\_v2} est la taille \textbf{recommandée par Microsoft} en
remplacement.

\textbf{4. Pourquoi ne pas conserver des IP publiques directes sur les VM en production ?}\\
Une IP publique par VM \textbf{expose chaque machine à Internet} (surface d'attaque, scans, brute-force),
\textbf{coûte en continu} et complique la gouvernance. En production, les VM restent en \textbf{IP privée} :
l'entrée se fait par un \textbf{Load Balancer / Application Gateway} et l'administration par \textbf{Azure
Bastion}. Ici, les IP publiques sont conservées pour \textbf{tester chaque VM individuellement}.

\begin{etatbox}
\begin{itemize}
  \item \code{templates/cloud-init.yml} + \code{compute.tf} créés : \textbf{2 IP publiques}, \textbf{2 interfaces réseau}, \textbf{2 VM Linux} \code{Standard\_B2ats\_v2} ; \code{apply} : \textbf{6 ressources ajoutées}.
  \item Nginx installé automatiquement par cloud-init ; pages distinctes (« serveur web 1 » / « serveur web 2 ») accessibles en HTTP.
  \item État Terraform synchronisé (\textbf{11 ressources} gérées).
  \item \textbf{Prêt pour l'Atelier 7 — mise en place du Load Balancer.}
\end{itemize}
\end{etatbox}

% #####################################################################
%  ATELIER 7
% #####################################################################
\section{Atelier 7 — Mise en place du Load Balancer}

\begin{fiche}
\textbf{\eTarget~Objectif} — répartir le trafic web entrant sur les deux VM via un Load Balancer public.

\textbf{\ePackage~Livrable attendu} — \code{loadbalancer.tf} et la preuve de répartition (alternance des réponses entre les 2 serveurs).

\textbf{\eGlobe~Région} — \code{swedencentral} \quad\textbullet\quad Load Balancer \code{lb-shopeasy-dev-web} (SKU Standard, public).
\end{fiche}

\subsection{IP publique, LB, pool, sonde et règle — \code{loadbalancer.tf}}
\begin{codebox}{hcl}
\begin{Verbatim}
resource "azurerm_public_ip" "lb" {
  name                = "pip-${local.prefix}-lb"
  location            = azurerm_resource_group.main.location
  resource_group_name = azurerm_resource_group.main.name
  allocation_method   = "Static"
  sku                 = "Standard"
  tags                = local.common_tags
}

resource "azurerm_lb" "web" {
  name                = "lb-${local.prefix}-web"
  location            = azurerm_resource_group.main.location
  resource_group_name = azurerm_resource_group.main.name
  sku                 = "Standard"
  tags                = local.common_tags
  frontend_ip_configuration {
    name                 = "public-frontend"
    public_ip_address_id = azurerm_public_ip.lb.id
  }
}

resource "azurerm_lb_backend_address_pool" "web" {
  name            = "backend-web"
  loadbalancer_id = azurerm_lb.web.id
}

resource "azurerm_network_interface_backend_address_pool_association" "web" {
  count                   = 2
  network_interface_id    = azurerm_network_interface.web[count.index].id
  ip_configuration_name   = "internal"
  backend_address_pool_id = azurerm_lb_backend_address_pool.web.id
}

resource "azurerm_lb_probe" "http" {
  name            = "http-probe"
  loadbalancer_id = azurerm_lb.web.id
  protocol        = "Http"
  request_path    = "/"
  port            = 80
}

resource "azurerm_lb_rule" "http" {
  name                           = "http-rule"
  loadbalancer_id                = azurerm_lb.web.id
  protocol                       = "Tcp"
  frontend_port                  = 80
  backend_port                   = 80
  frontend_ip_configuration_name = "public-frontend"
  backend_address_pool_ids       = [azurerm_lb_backend_address_pool.web.id]
  probe_id                       = azurerm_lb_probe.http.id
}
\end{Verbatim}
\end{codebox}

\begin{ltable}{X[l,0.9] X[l,1.1]}
Ressource & Rôle \\
\code{azurerm\_public\_ip.lb} & IP publique statique du frontend du LB. \\
\code{azurerm\_lb.web} & Load Balancer \textbf{Standard} (L4), frontend \code{public-frontend}. \\
\code{azurerm\_lb\_backend\_address\_pool.web} & Pool des VM cibles (\code{backend-web}). \\
\code{azurerm\_network\_interface\_}\allowbreak\code{backend\_address\_pool\_association.web} (×2) & Rattache chaque NIC (\code{web-1}, \code{web-2}) au pool. \\
\code{azurerm\_lb\_probe.http} & Sonde HTTP sur \code{/} port 80 (santé des VM). \\
\code{azurerm\_lb\_rule.http} & Règle 80 \arr{} 80 (Tcp) liant frontend + pool + sonde. \\
\end{ltable}

\subsection{Plan et application}
\begin{codebox}{console}
\begin{Verbatim}
Plan: 7 to add, 0 to change, 0 to destroy.
(LB + frontend, backend pool, 2 associations NIC, sonde, regle, IP publique)

azurerm_public_ip.lb: Creation complete after 2s [id=.../pip-shopeasy-dev-lb]
azurerm_lb.web: Creation complete after 11s [id=.../lb-shopeasy-dev-web]
azurerm_lb_probe.http: Creation complete after 11s
azurerm_lb_backend_address_pool.web: Creation complete after 14s
azurerm_network_interface_backend_address_pool_association.web[0..1]: complete after 2s
azurerm_lb_rule.http: Creation complete after 9s

Apply complete! Resources: 7 added, 0 changed, 0 destroyed.
\end{Verbatim}
\end{codebox}

\subsection{Vérification et test de répartition}
IP publique du Load Balancer : \texttt{9.223.33.133} (SKU Standard, frontend \code{public-frontend}).

\smallskip
\noindent\begin{minipage}[t]{0.48\linewidth}
\begin{ltable}{Q[l] Q[c] Q[c] Q[c]}
Règle & Proto & Front & Back \\
\code{http-rule} & Tcp & 80 & 80 \\
\end{ltable}
\end{minipage}\hfill
\begin{minipage}[t]{0.48\linewidth}
\begin{ltable}{Q[l] Q[c] Q[c] Q[c]}
Sonde & Proto & Port & Chemin \\
\code{http-probe} & Http & 80 & / \\
\end{ltable}
\end{minipage}

\smallskip
Les deux interfaces \code{nic-shopeasy-dev-web-1/2} sont membres du pool \code{backend-web}. Test de
répartition — \textbf{12 requêtes} sur l'IP du Load Balancer :
\begin{codebox}{console}
\begin{Verbatim}
$ for i in $(seq 1 12); do curl -s http://9.223.33.133 | grep -o "serveur web [12]"; done
req 1->web 2  req 2->web 1  req 3->web 1  req 4->web 1  req 5->web 2  req 6->web 1
req 7->web 2  req 8->web 2  req 9->web 2  req 10->web 1 req 11->web 1 req 12->web 2
=> 6 reponses "serveur web 1" / 6 reponses "serveur web 2"
\end{Verbatim}
\end{codebox}
\begin{tipbox}[Répartition de charge effective]
Sur 12 requêtes via le LB, \textbf{6 sur web-1 et 6 sur web-2} : le Load Balancer répartit bien le trafic
entre les deux VM, et la sonde HTTP confirme leur état de santé. C'est la suppression du point de défaillance
unique applicatif via un \textbf{point d'entrée unique}.
\end{tipbox}

\subsection{Captures}
\screenshot{tp2-lb-overview.png}{Vue d'ensemble du Load Balancer \code{lb-shopeasy-dev-web} (frontend, backend pool, règle) — \textit{subscription ID flouté}.}
\screenshot{tp2-web-lb.png}{Page web servie via le Load Balancer (\code{http://9.223.33.133}) ; en rafraîchissant, la réponse alterne entre « serveur web 1 » et « serveur web 2 ».}

\subsection{Analyse}
\textbf{1. Quel problème résout le Load Balancer ?}\\
Il \textbf{répartit le trafic entrant} sur plusieurs VM derrière un \textbf{point d'entrée public unique},
évitant la surcharge d'un seul serveur, améliorant la \textbf{disponibilité} et permettant de \textbf{monter en
charge}. Il découple l'adresse publique des serveurs individuels.

\textbf{2. Que se passe-t-il si une VM devient indisponible ?}\\
La \textbf{sonde de santé} (\code{http-probe} sur \code{/}) cesse de recevoir une réponse \code{200} et
\textbf{retire la VM de la rotation}. Le trafic est routé \textbf{uniquement vers la (les) VM saine(s)} : le
service reste disponible. Dès que la VM repasse la sonde, elle est \textbf{réintégrée automatiquement}.

\textbf{3. Quelle différence avec Azure Application Gateway ?}\\
\textbf{Azure Load Balancer} opère en \textbf{couche 4} (TCP/UDP) : répartition par hash, sans comprendre le
contenu HTTP. \textbf{Application Gateway} opère en \textbf{couche 7} (HTTP/HTTPS) : routage par chemin/hôte,
\textbf{terminaison TLS}, affinité de session et surtout \textbf{WAF}. On choisit le LB pour une répartition L4
générique, et l'Application Gateway pour les fonctions L7 (cf. analyse de sécurité).

\begin{etatbox}
\begin{itemize}
  \item \code{loadbalancer.tf} créé : IP publique + Load Balancer Standard + backend pool (2 VM) + sonde HTTP + règle 80 ; \code{apply} : \textbf{7 ressources ajoutées}.
  \item Test validé : \textbf{alternance 6/6} des réponses entre « serveur web 1 » et « serveur web 2 » via l'IP du LB.
  \item État Terraform : \textbf{18 ressources} gérées.
  \item \textbf{Prêt pour l'Atelier 8 — ajout d'un Storage Account.}
\end{itemize}
\end{etatbox}


% #####################################################################
%  ATELIER 8
% #####################################################################
\section{Atelier 8 — Ajout d'un Storage Account}

\begin{fiche}
\textbf{\eTarget~Objectif} — ajouter un stockage documentaire privé et versionné, avec un nom globalement unique.

\textbf{\ePackage~Livrable attendu} — \code{storage.tf} (suffixe aléatoire + Storage Account + container) et vérifications (privé, versioning, tags).
\end{fiche}

\subsection{Suffixe aléatoire, Storage Account et container — \code{storage.tf}}
\begin{codebox}{hcl}
\begin{Verbatim}
resource "random_string" "suffix" {
  length  = 6
  upper   = false
  special = false
}

resource "azurerm_storage_account" "docs" {
  name                     = "${replace(local.prefix, "-", "")}docs${random_string.suffix.result}"
  resource_group_name      = azurerm_resource_group.main.name
  location                 = azurerm_resource_group.main.location
  account_tier             = "Standard"
  account_replication_type = "LRS"
  min_tls_version          = "TLS1_2"
  tags                     = local.common_tags
  blob_properties {
    versioning_enabled = true
  }
}

resource "azurerm_storage_container" "documents" {
  name                  = "documents"
  storage_account_id    = azurerm_storage_account.docs.id
  container_access_type = "private"
}
\end{Verbatim}
\end{codebox}

\begin{ltable}{Q[l,0.5] X[l,1]}
Élément / choix & Raison \\
\code{random\_string.suffix} (6 car.) & Le nom d'un Storage Account doit être \textbf{globalement unique} (3-24 car., minuscules/chiffres). \\
\code{replace(local.prefix,"-","")} & Un Storage Account \textbf{n'accepte ni tiret ni majuscule} \arr{} \code{shopeasydevdocs1ac6q9}. \\
\code{account\_replication\_type = "LRS"} & Redondance locale, suffisante en dev (pas de GRS facturée inutilement). \\
\code{min\_tls\_version = "TLS1\_2"} & Refuse les protocoles obsolètes (sécurité en transit). \\
\code{versioning\_enabled = true} & Conserve l'historique des objets (récupération après écrasement/suppression). \\
\code{container\_access\_type = "private"} & Aucun accès anonyme : seuls les appels authentifiés accèdent aux documents. \\
\end{ltable}
Le provider \code{random} génère le suffixe \textbf{au moment de l'\code{apply}} : sa valeur est « known after
apply » dans le plan.

\subsection{Plan, application et vérification}
\begin{codebox}{console}
\begin{Verbatim}
Plan: 3 to add, 0 to change, 0 to destroy.   (random_string + Storage Account + container)
random_string.suffix: Creation complete [id=1ac6q9]
azurerm_storage_account.docs: Creation complete after 1m6s [id=.../shopeasydevdocs1ac6q9]
azurerm_storage_container.documents: Creation complete after 11s [id=.../containers/documents]
Apply complete! Resources: 3 added, 0 changed, 0 destroyed.

# Verifications Azure CLI
storage account : Standard_LRS, StorageV2, TLS1_2, httpsOnly=true
versioning      : isVersioningEnabled = true
container        : "documents" -> publicAccess = null  (= PRIVE)
\end{Verbatim}
\end{codebox}
Le suffixe généré est \code{1ac6q9} \arr{} Storage Account \textbf{\code{shopeasydevdocs1ac6q9}}. Bilan :
\code{StorageV2} / \code{Standard\_LRS}, \textbf{TLS 1.2}, HTTPS-only, \textbf{versioning activé}, et container
\code{documents} \textbf{privé}.

\subsection{Captures portail}
\screenshot{tp2-storage-overview.png}{Storage Account \code{shopeasydevdocs1ac6q9} (\code{Standard\_LRS}, \code{StorageV2}, tags) — \textit{subscription ID flouté}.}
\screenshot{tp2-storage-container.png}{Container \code{documents} avec un niveau d'accès \textit{Private}.}
\screenshot{tp2-storage-versioning.png}{Versioning Blob activé (Data protection).}

\subsection{Questions FinOps et sécurité}
\textbf{1. Pourquoi le container doit-il être privé ?}\\
Le container héberge des \textbf{documents métier sensibles} (factures, pièces clients). En accès \textbf{privé},
seuls les appels \textbf{authentifiés et autorisés} (clé de compte, jeton SAS, ou identité Azure AD via RBAC)
peuvent lire les blobs. Un container \textbf{public} exposerait tous les objets à quiconque connaît l'URL \arr{}
\textbf{fuite de données} et non-conformité (RGPD).

\textbf{2. Pourquoi le versioning peut-il augmenter les coûts ?}\\
À chaque modification ou suppression, le versioning \textbf{conserve une copie} de la version précédente. Ces
versions \textbf{s'accumulent} et sont \textbf{facturées comme du stockage}. Sans gestion du cycle de vie, les
anciennes versions ne sont \textbf{jamais nettoyées}.

\textbf{3. Quelle règle de cycle de vie proposer pour limiter les coûts ?}\\
Une \textbf{lifecycle management policy} : transférer les versions anciennes vers un niveau \textbf{Cool} puis
\textbf{Archive} après \textit{N} jours, et \textbf{supprimer} les versions de plus de \textit{X} jours. Cela
\textbf{plafonne} l'accumulation des versions et le coût associé.

\subsection{Note de sécurité — accès public au niveau du compte}
La vérification montre \code{allowBlobPublicAccess = true} \textbf{au niveau du compte} : le container est bien
privé, mais le compte \textbf{autoriserait} qu'un container soit passé en public. Pour la \textbf{défense en
profondeur}, on recommande de désactiver explicitement cette possibilité :
\begin{codebox}{hcl}
\begin{Verbatim}
# dans azurerm_storage_account.docs
allow_nested_items_to_be_public = false
\end{Verbatim}
\end{codebox}
Ainsi, aucun blob ne peut être rendu public, même par erreur — mesure cohérente avec le moindre privilège
(reprise dans l'analyse de sécurité de l'Atelier 13).

\begin{etatbox}
\begin{itemize}
  \item \code{storage.tf} créé : \code{random\_string} + Storage Account \code{shopeasydevdocs1ac6q9} (\code{Standard\_LRS}, \code{StorageV2}, TLS 1.2) + container \code{documents} \textbf{privé} \textbf{versionné} ; \code{apply} : \textbf{3 ressources ajoutées}.
  \item Risque résiduel identifié : \code{allowBlobPublicAccess = true} au niveau du compte (mesure corrective proposée).
  \item État Terraform : \textbf{21 ressources} gérées.
  \item \textbf{Prêt pour l'Atelier 9 — outputs et validation finale.}
\end{itemize}
\end{etatbox}

% #####################################################################
%  ATELIER 9
% #####################################################################
\section{Atelier 9 — Outputs et validation}

\begin{fiche}
\textbf{\eTarget~Objectif} — exposer les informations utiles après déploiement et valider l'ensemble de l'infrastructure.

\textbf{\ePackage~Livrable attendu} — \code{outputs.tf}, les sorties \code{terraform output}, la checklist technique et la page web via le Load Balancer.
\end{fiche}

\subsection{Sorties exposées — \code{outputs.tf}}
\begin{codebox}{hcl}
\begin{Verbatim}
output "resource_group_name" {
  description = "Nom du Resource Group cree"
  value       = azurerm_resource_group.main.name
}
output "load_balancer_public_ip" {
  description = "IP publique du Load Balancer (point d'entree de l'application)"
  value       = azurerm_public_ip.lb.ip_address
}
output "web_vm_public_ips" {
  description = "IP publiques des deux VM web"
  value       = azurerm_public_ip.web[*].ip_address
}
output "storage_account_name" {
  description = "Nom du Storage Account documentaire"
  value       = azurerm_storage_account.docs.name
}
\end{Verbatim}
\end{codebox}
Les outputs exposent les informations dont un développeur ou un script a besoin \textbf{après} le déploiement,
sans fouiller le portail. L'expression \code{azurerm\_public\_ip.web[*].ip\_address} utilise le \textit{splat}
\code{[*]} pour renvoyer la \textbf{liste} des IP des deux VM.

\subsection{Application et lecture des sorties}
Ajouter des outputs \textbf{ne modifie aucune ressource} : Terraform met seulement à jour les sorties du state.
\begin{codebox}{console}
\begin{Verbatim}
Changes to Outputs:
  + load_balancer_public_ip = "9.223.33.133"
  + resource_group_name     = "rg-shopeasy-dev"
  + storage_account_name    = "shopeasydevdocs1ac6q9"
  + web_vm_public_ips        = [ "135.225.73.252", "135.225.74.4" ]

Apply complete! Resources: 0 added, 0 changed, 0 destroyed.

$ terraform output
load_balancer_public_ip = "9.223.33.133"
resource_group_name     = "rg-shopeasy-dev"
storage_account_name    = "shopeasydevdocs1ac6q9"
web_vm_public_ips       = [ "135.225.73.252", "135.225.74.4" ]
\end{Verbatim}
\end{codebox}
{\small \code{terraform show} a également été exécuté : il restitue l'intégralité du state lisible (les 21 ressources gérées et leurs attributs). Sortie volumineuse, non reproduite ici.}

\subsection{Checklist technique}
\begin{ltable}{X[l,1.1] Q[c,0.2] X[l,1.3]}
Contrôle & État & Preuve \\
Le projet Terraform s'initialise sans erreur & \eCheck & Atelier 2 — \code{terraform init} : providers installés. \\
\code{terraform validate} réussit & \eCheck & \textit{Success! The configuration is valid.} (chaque atelier). \\
Le Resource Group est créé & \eCheck & Atelier 4 — \code{rg-shopeasy-dev}, \code{Succeeded}. \\
Le VNet et le subnet existent & \eCheck & Atelier 4 — \code{vnet-shopeasy-dev} + \code{snet-web}. \\
Le NSG limite SSH à l'IP admin & \eCheck & Atelier 5 — \code{Allow-SSH-Admin} 22 depuis \code{<IP\_ADMIN>/32}. \\
Deux VM Linux sont créées & \eCheck & Atelier 6 — \code{vm-...-web-1/2}, \code{VM running}. \\
Nginx répond sur les VM & \eCheck & Atelier 6 — \code{curl} \arr{} « serveur web 1 » / « 2 ». \\
Le Load Balancer répond en HTTP & \eCheck & Atelier 7 — alternance \textbf{6/6} sur \code{http://9.223.33.133}. \\
Le Storage Account est privé & \eCheck & Atelier 8 — container \code{documents} \code{publicAccess: null}. \\
Le versioning Blob est activé & \eCheck & Atelier 8 — \code{isVersioningEnabled: true}. \\
Les ressources sont taguées & \eCheck & \code{local.common\_tags} (4 tags) sur RG, VNet, NSG, IP, NIC, VM, LB, Storage. \\
\end{ltable}
Tous les contrôles sont \textbf{au vert} (11/11).

\subsection{Page web via le Load Balancer}
\screenshot{tp2-web-lb.png}{Page web accessible via le Load Balancer (\code{http://9.223.33.133}) : en rafraîchissant, la réponse alterne entre les deux serveurs (répartition de charge).}

\begin{etatbox}
\begin{itemize}
  \item \code{outputs.tf} créé : \textbf{4 sorties} (nom du RG, IP du LB, IP des VM, nom du Storage).
  \item \code{terraform apply} : \textbf{0 ressource modifiée}, sorties ajoutées au state ; checklist technique \textbf{11/11 OK}.
  \item Infrastructure complète : \textbf{21 ressources} gérées, déployées et validées.
  \item \textbf{Prêt pour l'Atelier 10 — modification de l'infrastructure.}
\end{itemize}
\end{etatbox}


% #####################################################################
%  ATELIER 10
% #####################################################################
\section{Atelier 10 — Modification de l'infrastructure}

\begin{fiche}
\textbf{\eTarget~Objectif} — modifier l'infrastructure existante (ajout d'un tag) et lire le plan pour comprendre comment Terraform applique un changement.

\textbf{\ePackage~Livrable attendu} — modification de \code{locals.tf} et analyse du plan (mise à jour vs recréation).
\end{fiche}

\subsection{Ajout du tag \code{cost\_center} dans \code{locals.tf}}
\begin{codebox}{hcl}
\begin{Verbatim}
common_tags = {
  project     = var.project
  environment = var.environment
  owner       = "formation"
  managed_by  = "terraform"
  cost_center = "cloud-training"   # <- tag ajoute
}
\end{Verbatim}
\end{codebox}
Comme ce \textit{local} alimente l'attribut \code{tags} de toutes les ressources concernées, \textbf{une seule
modification} se propage à l'ensemble du projet (principe DRY).

\subsection{Prévisualisation — \code{terraform plan}}
\begin{codebox}{console}
\begin{Verbatim}
  # azurerm_resource_group.main will be updated in-place
  ~ resource "azurerm_resource_group" "main" {
        name = "rg-shopeasy-dev"
      ~ tags = {
          + "cost_center" = "cloud-training"
            "environment" = "dev"
            "managed_by"  = "terraform"
            "owner"       = "formation"
            "project"     = "shopeasy"
        }
    }
  # ... (11 autres ressources taguees : VNet, NSG, IP, NIC, VM, LB, Storage)

Plan: 0 to add, 12 to change, 0 to destroy.
\end{Verbatim}
\end{codebox}
Le \code{\textasciitilde} indique une \textbf{mise à jour en place} et le \code{+} l'\textbf{ajout} d'une seule
clé de tag ; les autres tags restent inchangés. Aucune ligne \code{-} (destruction) ni \code{-/+}
(remplacement) n'apparaît.

\subsection{Application et idempotence}
\begin{codebox}{console}
\begin{Verbatim}
azurerm_resource_group.main: Modifications complete after 1s
...
Apply complete! Resources: 0 added, 12 changed, 0 destroyed.

# Nouveau plan (idempotence)
Terraform has compared your real infrastructure against your configuration
and found no differences, so no changes are needed.

# Tag propage (verification)
VNet : cloud-training   VM-1 : cloud-training   LB : cloud-training   Store : cloud-training
\end{Verbatim}
\end{codebox}

\subsection{Analyse du plan}
\textbf{1. Terraform prévoit-il de recréer toutes les ressources ?}\\
\textbf{Non.} Le plan annonce \code{0 to add, 12 to change, 0 to destroy} et chaque ressource est marquée
\textbf{\code{update in-place}} (\code{\textasciitilde}). Un tag est un attribut \textbf{modifiable à chaud} :
Terraform l'ajoute directement, \textbf{sans interruption de service}.

\textbf{2. Quelles ressources sont simplement mises à jour ?}\\
Les \textbf{12 ressources qui portent l'attribut \code{tags}} (= \code{local.common\_tags}) : RG, VNet, NSG,
les 2 IP web + l'IP du LB, les 2 NIC, les 2 VM, le LB et le Storage. Les ressources \textbf{sans tags} (subnet,
associations, backend pool, sonde, règle, container, \code{random\_string}) ne sont \textbf{pas concernées}.

\textbf{3. Pourquoi le plan est-il indispensable avant application ?}\\
Le plan montre \textbf{exactement} ce qui va changer \textbf{avant} d'agir : ici, des tags ajoutés en place
(\code{\textasciitilde}), sans destruction (\code{-}) ni remplacement (\code{-/+}). Lire le plan permet de
\textbf{détecter une action dangereuse} (destruction imprévue, remplacement de VM effaçant des données) avant
qu'elle ne se produise. C'est le garde-fou du caractère \textbf{déclaratif} de Terraform.

\begin{etatbox}
\begin{itemize}
  \item \code{locals.tf} modifié : ajout du tag \code{cost\_center = "cloud-training"} ; \code{apply} : \textbf{12 ressources mises à jour en place}, 0 création, 0 destruction.
  \item Idempotence vérifiée (\code{plan} suivant : \textit{no changes}), tag propagé à toutes les ressources taguées.
  \item Démonstration de la différence \textbf{mise à jour (\code{\textasciitilde}) vs recréation (\code{-/+})}.
  \item \textbf{Prêt pour l'Atelier 11 — observation d'une dérive (drift).}
\end{itemize}
\end{etatbox}


% #####################################################################
%  ATELIER 11
% #####################################################################
\section{Atelier 11 — Observation d'une dérive Terraform}

\begin{fiche}
\textbf{\eTarget~Objectif} — provoquer une modification manuelle hors Terraform et observer comment l'outil détecte l'écart (\textit{drift}).

\textbf{\ePackage~Livrable attendu} — mise en évidence de la dérive via \code{terraform plan} et analyse.
\end{fiche}

\subsection{Création d'une dérive volontaire}
Une modification est faite \textbf{manuellement dans le portail Azure}, hors Terraform : ajout du tag
\code{manual\_change = true} sur le Resource Group \code{rg-shopeasy-dev}. Le RG porte alors \textbf{6 tags}
(les 5 déclarés + le tag manuel). Cette modification \textbf{n'existe pas dans le code} : c'est précisément ce
qui constitue une dérive.

\subsection{Détection — \code{terraform plan}}
\begin{codebox}{console}
\begin{Verbatim}
  # azurerm_resource_group.main will be updated in-place
  ~ resource "azurerm_resource_group" "main" {
        name = "rg-shopeasy-dev"
      ~ tags = {
            "cost_center"   = "cloud-training"
            "environment"   = "dev"
            "managed_by"    = "terraform"
          - "manual_change" = "true" -> null
            "owner"         = "formation"
            "project"       = "shopeasy"
        }
    }

Plan: 0 to add, 1 to change, 0 to destroy.
\end{Verbatim}
\end{codebox}
Terraform \textbf{compare l'état réel au code} (sa source de vérité) et détecte le tag \code{manual\_change}
ajouté hors processus. Il propose de le \textbf{supprimer} (\code{- "manual\_change" = "true" -> null}) pour
\textbf{réaligner la réalité sur l'état déclaré}.

\subsection{Réconciliation — \code{terraform apply}}
\begin{codebox}{console}
\begin{Verbatim}
azurerm_resource_group.main: Modifications complete after 1s
Apply complete! Resources: 0 added, 1 changed, 0 destroyed.

# Tags apres reconciliation : retour aux 5 tags declares (manual_change supprime)
# Plan de controle
Terraform has compared your real infrastructure against your configuration
and found no differences, so no changes are needed.
\end{Verbatim}
\end{codebox}
La dérive est résorbée : l'infrastructure correspond de nouveau exactement au code.

\subsection{Capture portail}
\screenshot{tp2-drift-manual.png}{Tag \code{manual\_change = true} ajouté manuellement sur le Resource Group (dérive volontaire, hors Terraform).}

\subsection{Analyse}
\textbf{1. Terraform détecte-t-il une différence ?}\\
\textbf{Oui.} Lors du \code{plan}, Terraform rafraîchit l'état réel, le compare au code et \textbf{signale
l'écart} : le tag \code{manual\_change} présent dans Azure n'existe pas dans la configuration.

\textbf{2. Quelle action propose-t-il ?}\\
Une \textbf{mise à jour en place} (\code{\textasciitilde}) du Resource Group pour \textbf{supprimer} le tag
\code{manual\_change}, afin de revenir à l'état déclaré (\code{0 to add, 1 to change, 0 to destroy}).

\textbf{3. Pourquoi les modifications manuelles sont-elles dangereuses dans une organisation ?}\\
Elles créent une \textbf{dérive} entre le code (source de vérité) et l'infrastructure réelle : environnement
\textbf{imprévisible}, \textbf{sécurité affaiblie} (une règle ouverte à la main échappe à la revue),
\textbf{confiance dans le code} perdue, et un \code{apply} ultérieur peut \textbf{annuler silencieusement} un
changement manuel — ou l'inverse. L'environnement devient difficile à auditer et reproduire.

\textbf{4. Quelle règle d'équipe proposer pour limiter ce risque ?}\\
Tout changement doit \textbf{passer par le code} : modification du \code{.tf} \arr{} \textbf{pull request} \arr{}
revue \arr{} \code{plan} \arr{} \code{apply} contrôlé. On complète par : \code{plan} régulier en CI (détection
des dérives), \textbf{droits portail restreints} (lecture seule en prod), \textbf{Azure Policy}, et formation
des équipes. Le portail sert alors à \textbf{observer et diagnostiquer}, pas à modifier.

\begin{etatbox}
\begin{itemize}
  \item Dérive volontaire créée (tag \code{manual\_change} ajouté manuellement) puis \textbf{détectée} par \code{terraform plan} (\code{0 add, 1 change, 0 destroy}).
  \item \textbf{Réconciliation} appliquée : retour aux 5 tags déclarés, \code{plan} de contrôle \textit{no changes}.
  \item Démonstration du rôle du code comme \textbf{source de vérité} et de l'intérêt du \code{plan}.
  \item \textbf{Prêt pour l'Atelier 12 — préparation d'un state distant.}
\end{itemize}
\end{etatbox}

% #####################################################################
%  ATELIER 12
% #####################################################################
\section{Atelier 12 — Préparation d'un state distant}

\begin{fiche}
\textbf{\eTarget~Objectif} — comprendre pourquoi et comment externaliser le state Terraform dans un backend distant sécurisé.

\textbf{\ePackage~Livrable attendu} — structure cible du backend Azure et analyse (atelier de \textbf{préparation conceptuelle}, sans migration du state).
\end{fiche}

\subsection{Le state local et ses limites}
Le projet utilise un \textbf{state local} : un fichier \code{terraform.tfstate} sur le poste. Ce state
\textbf{contient des données sensibles en clair} — une recherche des \textbf{noms} d'attributs sensibles (sans
afficher aucune valeur) le confirme :
\begin{codebox}{console}
\begin{Verbatim}
$ ls -la terraform.tfstate          # 59 639 octets
$ git check-ignore terraform.tfstate   # -> ignore (non versionne)
$ terraform state list | wc -l         # 21 ressources suivies

# Recherche de noms d'attributs sensibles dans le state :
cle presente: "admin_password"
cle presente: "primary_access_key"     <- cle d'acces du Storage Account
cle presente: "secret"
\end{Verbatim}
\end{codebox}
Le state local est acceptable en découverte individuelle, mais \textbf{insuffisant en équipe} :
\begin{itemize}
  \item \textbf{perte} possible du fichier ; \textbf{conflits} entre plusieurs administrateurs ;
  \item \textbf{pas de verrouillage} \arr{} corruption en cas d'\code{apply} simultané ;
  \item \textbf{présence de secrets} \arr{} un commit accidentel exposerait clés et mots de passe ;
  \item \textbf{difficile à intégrer} dans une chaîne CI/CD.
\end{itemize}

\subsection{Structure cible — backend Azure Storage}
\begin{codebox}{hcl}
\begin{Verbatim}
terraform {
  backend "azurerm" {
    resource_group_name  = "rg-terraform-state"
    storage_account_name = "stterraformstatexxx"
    container_name       = "tfstate"
    key                  = "shopeasy/dev/terraform.tfstate"
  }
}
\end{Verbatim}
\end{codebox}
\begin{ltable}{Q[l,0.6] X[l,1]}
Champ & Rôle \\
\code{resource\_group\_name} & RG dédié au state (séparé des ressources applicatives). \\
\code{storage\_account\_name} & Storage Account hébergeant le state (nom globalement unique). \\
\code{container\_name} & Container Blob \code{tfstate}. \\
\code{key} & Chemin/nom du blob de state — \textbf{inclut l'environnement} (\code{shopeasy/dev/...}) pour isoler dev/recette/prod. \\
\end{ltable}
Le backend \code{azurerm} fournit nativement le \textbf{verrouillage} (via un \textit{lease} sur le blob) : deux
\code{apply} simultanés ne peuvent pas corrompre le state.

\begin{notebox}[Bootstrap — ne pas ajouter le bloc à la légère]
Le bloc \code{backend} ne doit \textbf{pas} être ajouté sans avoir créé \textbf{au préalable} le Resource
Group, le Storage Account et le container de state. Dans un vrai projet, cette initialisation est réalisée par
une \textbf{équipe plateforme} ou un \textbf{bootstrap dédié}, puis sécurisée (RBAC, chiffrement, réseau). La
bascule se ferait ensuite avec \code{terraform init -migrate-state}. \textbf{Cette migration n'est pas
effectuée ici} : l'atelier reste conceptuel, et le state du TP demeure local (puis supprimé à l'Atelier 14).
\end{notebox}

\subsection{Questions}
\textbf{1. Pourquoi le state distant facilite-t-il le travail en équipe ?}\\
Il fournit un \textbf{emplacement central et partagé} : tous les membres et la \textbf{CI/CD} lisent et écrivent
le \textbf{même} state, sans copies divergentes. Il apporte le \textbf{verrouillage} (un seul \code{apply} à la
fois), supprime le state des postes (\textbf{pas de perte ni de commit accidentel}) et s'intègre dans un
\textbf{pipeline}.

\textbf{2. Pourquoi faut-il protéger l'accès au Storage Account du state ?}\\
Parce que le state \textbf{contient des secrets en clair} (\code{admin\_password}, \code{primary\_access\_key},
\code{secret}) \textbf{et} l'inventaire complet de l'infrastructure. Il faut : \textbf{RBAC restreint},
\textbf{chiffrement au repos}, \textbf{restrictions réseau} (private endpoint / pare-feu), désactivation de
l'accès public, rotation des clés et journalisation.

\textbf{3. Pourquoi faut-il séparer le state de développement, recette et production ?}\\
Pour \textbf{isoler les environnements} : une erreur sur le state \textbf{dev} ne doit jamais impacter la
\textbf{prod}. La séparation (clé/container distincts, ex. \code{shopeasy/prod/terraform.tfstate}) permet des
\textbf{droits différenciés}, évite d'\textbf{appliquer par erreur} un changement dev sur la prod, et donne à
chaque environnement un \textbf{cycle de vie indépendant}.

\begin{etatbox}
\begin{itemize}
  \item Limites du state local démontrées : fichier local de 21 ressources \textbf{contenant des secrets en clair}.
  \item Structure cible documentée : backend \code{azurerm} (Storage Account + container \code{tfstate}, clé par environnement, verrouillage natif) ; bootstrap et migration expliqués (\textbf{non appliqués}).
  \item \textbf{Prêt pour l'Atelier 13 — analyses coût, sécurité et maintenabilité.}
\end{itemize}
\end{etatbox}


% #####################################################################
%  ATELIER 13
% #####################################################################
\section{Atelier 13 — Analyse coût, sécurité et maintenabilité}

\begin{fiche}
\textbf{\eTarget~Objectif} — analyser l'infrastructure déployée sous trois angles : coûts (FinOps), sécurité et maintenabilité.

\textbf{\ePackage~Livrable attendu} — tableau FinOps, tableau d'analyse de sécurité et réponses sur la maintenabilité.
\end{fiche}

\subsection{Analyse FinOps}
Prix unitaires \textbf{réels} (Azure Retail Prices API, région \code{swedencentral}, pay-as-you-go USD,
hypothèse 24/7 = 730 h/mois) :
\begin{ltable}{X[l,1.2] Q[l,0.7]}
Poste & Prix unitaire officiel \\
VM \code{Standard\_B2ats\_v2} (Linux) & \textbf{0,00972 \$/h} \\
IP publique Standard (static) & \textbf{0,005 \$/h} \\
Disque OS Standard HDD S4 (32 Go) LRS & \textbf{1,536 \$/mois} \\
Load Balancer Standard & \textbf{$\approx$ 0,025 \$/h} (règles incluses) \\
Blob Storage (faible volume) & négligeable \\
\end{ltable}

\begin{tblr}{
  width=\linewidth, colspec={Q[l,0.7] Q[l,0.6] X[l,1.1] X[l,1.3]},
  row{1}={bg=azuredark,fg=white,font=\sffamily\bfseries\scriptsize},
  row{2-Z}={font=\scriptsize}, row{even}={bg=rowalt},
  hline{1,Z}={1.1pt,azuredark}, hline{2}={0.5pt,azuredark!45},
  colsep=4pt, rowsep=3.2pt, stretch=1.15,
}
Ressource & Coût relatif & Risque de surcoût & Optimisation proposée \\
\textbf{VM Linux} (2 × \code{B2ats\_v2}) & Moyen-élevé ($\approx$ 14,2 \$) & VM allumées 24/7 sans usage réel & Désallouer hors période ; autoscale en prod ; taille adaptée \\
\textbf{IP publiques} (3) & Moyen ($\approx$ 10,9 \$) & Facturées même à trafic nul & Réduire à 1 (celle du LB) ; VM sans IP publique + Bastion \\
\textbf{Load Balancer} (Standard) & Élevé ($\approx$ 18,2 \$) & Facturé en continu, même sans trafic & Garder pour la prod ; App Gateway si L7/WAF requis \\
\textbf{Storage Account} (LRS) & Faible (< 0,1 \$) & Réplication géo (GRS) inutile en dev & LRS (pas GRS) ; suppression au nettoyage \\
\textbf{Versioning Blob} & Faible \arr{} croissant & Accumulation illimitée d'anciennes versions & Lifecycle : transition Cool/Archive + purge des versions \\
\textbf{Disques managés} (2 × OS, S4) & Faible ($\approx$ 3,1 \$) & Disques orphelins après suppression VM & Taille adaptée ; suppression au nettoyage \\
\end{tblr}

\begin{notebox}[Total et postes principaux]
\textbf{Total indicatif $\approx$ 46,5 \$/mois en 24/7.} Les 3 postes principaux sont le \textbf{Load Balancer
(18,2 \$)}, le \textbf{compute des 2 VM (14,2 \$)} et les \textbf{IP publiques (10,9 \$)} — le LB et les IP
facturent même à trafic nul. Désallouer les VM annule leur coût compute, et la \textbf{suppression du Resource
Group} (Atelier 14) ramène le coût à $\approx$ 0.
\end{notebox}

\subsection{Analyse de sécurité}
\begin{tblr}{
  width=\linewidth, colspec={Q[l,0.6] X[l,0.9] X[l,0.9] X[l,1.1]},
  row{1}={bg=azuredark,fg=white,font=\sffamily\bfseries\scriptsize},
  row{2-Z}={font=\scriptsize}, row{even}={bg=rowalt},
  hline{1,Z}={1.1pt,azuredark}, hline{2}={0.5pt,azuredark!45},
  colsep=4pt, rowsep=3.2pt, stretch=1.15,
}
Risque & Cause possible & Impact & Correction \\
\textbf{SSH trop ouvert} & Source \code{0.0.0.0/0} au lieu de l'IP admin & Scans, brute-force, compromission & Restreindre à l'IP \code{/32} \textit{(appliqué)} ; Bastion / JIT en prod \\
\textbf{State exposé} & \code{tfstate} commité ou Storage de state non protégé & Fuite de secrets + inventaire du SI & \code{.gitignore} \textit{(appliqué)} ; backend distant chiffré + RBAC \\
\textbf{Storage public} & Container public ou \code{allowBlobPublicAccess=true} & Fuite de documents métier (RGPD) & Container privé \textit{(appliqué)} ; \code{allow\_nested\_items\_to\_be\_public=false} \\
\textbf{Tags absents} & Oubli de tags sur des ressources & Coûts non imputables, gouvernance impossible & \code{common\_tags} systématiques \textit{(appliqué)} ; Azure Policy \\
\textbf{Secrets dans Git} & Mot de passe/clé en dur dans \code{.tf} ou \code{tfvars} versionné & Compromission d'identifiants & Aucun secret en dur ; \code{tfvars} ignoré \textit{(appliqué)} ; Key Vault \\
\textbf{Modification manuelle} & Changement portail hors Terraform & Dérive, imprévisibilité, sécurité affaiblie & Tout par le code (PR + plan + apply) ; \code{plan} régulier en CI \\
\end{tblr}
{\small Les corrections \textit{(appliqué)} sont déjà en place (Ateliers 1, 5, 8, 11). Les autres sont les évolutions recommandées pour la production.}

\subsection{Maintenabilité}
\textbf{1. Comment rendre ce projet réutilisable pour un environnement de recette ?}\\
Le code est déjà \textbf{paramétré par variables} : il suffit de fournir un \textbf{\code{terraform.tfvars} par
environnement} (ex. \code{environment = "recette"}, sa région, sa taille de VM) et un \textbf{state séparé}
(clé/backend distincts). Le \textbf{code reste unique} ; seules les \textbf{valeurs} changent. On peut aussi
utiliser des \textbf{workspaces}.

\textbf{2. Quelles variables faudrait-il ajouter ?}\\
\code{vm\_size}, \code{vm\_count} (nombre d'instances), \code{address\_space} / \code{web\_subnet\_prefix},
\code{storage\_replication\_type}, \code{enable\_public\_ip} (booléen pour supprimer les IP publiques en prod),
et un bloc de \textbf{tags additionnels}.

\textbf{3. Quels fichiers pourraient devenir des modules Terraform ?}\\
\code{network.tf} + \code{security.tf} \arr{} \textbf{module \code{network}} ; \code{compute.tf} \arr{}
\textbf{module \code{compute}} ; \code{loadbalancer.tf} \arr{} \textbf{module \code{loadbalancer}} ;
\code{storage.tf} \arr{} \textbf{module \code{storage}}. Appelés depuis un \code{main.tf} avec des variables,
ils deviennent \textbf{réutilisables} pour dev / recette / prod.

\textbf{4. Quelles validations automatiques pourrait-on ajouter dans une pipeline CI/CD ?}\\
\code{terraform fmt -check} et \code{validate} ; \textbf{\code{tflint}} ; \textbf{\code{tfsec} / \code{checkov}}
(sécurité statique) ; \textbf{\code{gitleaks}} (secrets) ; \code{terraform plan} automatique sur chaque
\textbf{pull request} + revue et approbation avant \code{apply} ; exécution régulière de \code{plan} pour
détecter les dérives.

\begin{etatbox}
\begin{itemize}
  \item \textbf{FinOps} : coût $\approx$ \textbf{46,5 \$/mois} (24/7, prix réels), 3 postes principaux (LB, compute, IP), optimisations par ressource.
  \item \textbf{Sécurité} : 6 risques analysés (cause / impact / correction), dont 4 déjà corrigés dans le projet.
  \item \textbf{Maintenabilité} : réutilisation par environnement, variables à ajouter, modularisation, validations CI/CD.
  \item \textbf{Prêt pour l'Atelier 14 — nettoyage de l'environnement.}
\end{itemize}
\end{etatbox}

% #####################################################################
%  ATELIER 14
% #####################################################################
\section{Atelier 14 — Nettoyage de l'environnement}

\begin{fiche}
\textbf{\eTarget~Objectif} — supprimer proprement toutes les ressources créées pour arrêter la facturation.

\textbf{\ePackage~Livrable attendu} — prévisualisation de la destruction et commande \code{terraform destroy}.
\end{fiche}

\begin{etatbox}[Destruction exécutée]
Après validation de l'ensemble des livrables et des captures, la destruction effective a été lancée
(\code{terraform destroy}). Les \textbf{21 ressources ont été supprimées} et le Resource Group n'existe plus
côté Azure — la \textbf{facturation est arrêtée} (détaillée ci-dessous).
\end{etatbox}

\subsection{Pourquoi détruire l'environnement ?}
L'infrastructure de test coûte $\approx$ \textbf{46,5 \$/mois en 24/7} (cf. Atelier 13), facturée \textbf{même
à trafic nul} pour le Load Balancer et les IP publiques. Comme tout est décrit en \textbf{Infrastructure as
Code}, l'environnement est \textbf{reproductible à l'identique} par \code{terraform apply} : il n'y a aucune
raison de le laisser tourner entre deux séances. La destruction ramène le coût à $\approx$ 0.

\subsection{Prévisualisation — \code{terraform plan -destroy}}
\begin{codebox}{console}
\begin{Verbatim}
$ terraform plan -destroy        # ne supprime rien : liste ce qui SERAIT detruit
azurerm_lb.web, azurerm_lb_backend_address_pool.web, azurerm_lb_probe.http,
azurerm_lb_rule.http, azurerm_linux_virtual_machine.web[0..1],
azurerm_network_interface.web[0..1],
azurerm_network_interface_backend_address_pool_association.web[0..1],
azurerm_network_security_group.web, azurerm_public_ip.lb,
azurerm_public_ip.web[0..1], azurerm_resource_group.main,
azurerm_storage_account.docs, azurerm_storage_container.documents,
azurerm_subnet.web, azurerm_subnet_network_security_group_association.web,
azurerm_virtual_network.main, random_string.suffix          ... will be destroyed

Plan: 0 to add, 0 to change, 21 to destroy.
\end{Verbatim}
\end{codebox}
Les \textbf{21 ressources} gérées par Terraform sont prévues pour suppression.

\subsection{Exécution et vérification}
Après validation de tous les livrables et captures, la destruction a été exécutée (\code{terraform destroy
-auto-approve}) : elle supprime toutes les ressources gérées par le state, dans l'ordre inverse des
dépendances. Sortie réelle (extrait, identifiant d'abonnement masqué) :
\begin{codebox}{console}
\begin{Verbatim}
random_string.suffix: Destruction complete after 0s
azurerm_lb_rule.http: Destruction complete after 4s
azurerm_network_security_group.web: Destruction complete after 11s
azurerm_lb.web: Destruction complete after 10s
azurerm_linux_virtual_machine.web[0]: Destruction complete after 32s
azurerm_linux_virtual_machine.web[1]: Destruction complete after 32s
azurerm_public_ip.lb / web[0..1]: Destruction complete after ~11s
azurerm_network_interface.web[0..1]: Destruction complete after ~11-22s
azurerm_subnet.web: Destruction complete after 11s
azurerm_virtual_network.main: Destruction complete after 11s
azurerm_resource_group.main: Destruction complete after 21s

Destroy complete! Resources: 21 destroyed.
\end{Verbatim}
\end{codebox}
Vérifications après destruction :
\begin{codebox}{console}
\begin{Verbatim}
$ terraform state list
(vide -- aucune ressource geree)

$ az group show -n rg-shopeasy-dev
ERROR: (ResourceGroupNotFound) Resource group 'rg-shopeasy-dev' could not be found.
\end{Verbatim}
\end{codebox}

\begin{notebox}[Point obligatoire du TP]
Ne jamais terminer une séance sans avoir supprimé les ressources de formation. C'est fait : les \textbf{21
ressources sont détruites}, le Resource Group n'existe plus, la \textbf{facturation est arrêtée}.
\end{notebox}

\begin{etatbox}
\begin{itemize}
  \item \code{terraform destroy} exécuté : \textbf{Destroy complete! Resources: 21 destroyed}.
  \item État Terraform \textbf{vide} ; Resource Group \code{rg-shopeasy-dev} \textbf{absent} de l'abonnement (\code{ResourceGroupNotFound}).
  \item \textbf{Facturation arrêtée} ($\approx$ 0) ; l'environnement reste \textbf{reproductible à l'identique} par \code{terraform apply}.
  \item \textbf{Fin du TP2} : projet Terraform complet — déployé, validé, documenté et nettoyé.
\end{itemize}
\end{etatbox}


% #####################################################################
%  NOTE TECHNIQUE
% #####################################################################
\section{Note technique — synthèse des choix}

\begin{tcolorbox}[enhanced, sharp corners, boxrule=0pt, leftrule=4pt, colback=azurepale, colframe=azuredark,
   left=4.5mm,right=4.5mm,top=3.6mm,bottom=3.6mm, before skip=10pt, after skip=12pt]
\textbf{Objet} — industrialisation du déploiement Azure de ShopEasy avec Terraform (Infrastructure as Code)\\[2pt]
\textbf{Portée} — synthèse des choix techniques, de sécurité, de coût et de gouvernance, et proposition d'évolution\\[2pt]
\textbf{Auteurs (binôme)} — Louis \textsc{Scarfone} \& Maxence \textsc{Bourragué}
\end{tcolorbox}

\subsection{Contexte}
Le TP1 avait conçu et déployé \textbf{manuellement} (Azure CLI) l'architecture cible de ShopEasy. Le TP2
reprend cette architecture et la décrit en \textbf{Infrastructure as Code} : l'infrastructure devient un
ensemble de fichiers \code{.tf} \textbf{versionnés, relus, appliqués puis détruits} de manière reproductible.
Le résultat n'est pas seulement « une infrastructure qui fonctionne », mais une infrastructure \textbf{décrite
proprement, lisible, versionnable, sécurisée et maîtrisée dans son cycle de vie}.

\subsection{Architecture déployée}
\begin{ltable}{Q[l,0.5] X[l,1.4]}
Couche & Ressources \\
\textbf{Organisation} & Resource Group \code{rg-shopeasy-dev} (\code{swedencentral}), tags de gouvernance. \\
\textbf{Réseau} & VNet \code{10.20.0.0/16}, subnet \code{snet-web} \code{10.20.1.0/24}, NSG (HTTP ouvert, SSH restreint à l'IP admin). \\
\textbf{Calcul} & 2 VM Linux \code{Standard\_B2ats\_v2} (Ubuntu 22.04), Nginx installé par cloud-init, déployées via \code{count}. \\
\textbf{Répartition} & Load Balancer \textbf{Standard} (IP publique, backend pool, sonde HTTP, règle 80). \\
\textbf{Stockage} & Storage Account \code{StorageV2} \textbf{LRS}, container \textbf{privé}, \textbf{versioning} Blob activé. \\
\textbf{Exploitation} & Outputs (IP du LB, IP des VM, nom du Storage), state local protégé. \\
\end{ltable}
\textbf{21 ressources} Terraform, déployées et validées (checklist technique de l'Atelier 9 : 11/11).

\subsection{Choix techniques}
\begin{ltable}{Q[l,0.55] Q[l,0.9] X[l,1.2]}
Choix & Décision & Justification \\
Outil IaC & \textbf{Terraform} + \code{azurerm \textasciitilde> 4.0} & Déclaratif, multi-cloud, lisible. \\
Authentification & \code{az login} + \code{ARM\_SUBSCRIPTION\_ID} & Aucun identifiant en dur dans le code. \\
Région & \textbf{\code{swedencentral}} & \code{francecentral} bloquée (policy Azure for Students). \\
Taille de VM & \textbf{\code{Standard\_B2ats\_v2}} & \code{B1s} indisponible ; \code{B2ats\_v2} recommandée Microsoft. \\
Multiplication & \code{count = 2} pour IP/NIC/VM & Deux serveurs identiques sans duplication (DRY). \\
Provisionnement & cloud-init + \code{templatefile} + \code{custom\_data} & Installation Nginx automatique, page par serveur. \\
Clé SSH & \code{file(pathexpand(...))} & \code{file()} ne développe pas \code{\textasciitilde}. \\
Nommage / tags & \code{local.prefix} + \code{local.common\_tags} & Source unique, cohérence et gouvernance. \\
Lisibilité & Découpage par responsabilité (\code{network}, \code{compute}…) & Projet relisible et maintenable. \\
\end{ltable}

\subsection{Sécurité}
SSH \textbf{restreint} à l'IP admin (\code{/32}) ; \textbf{aucun secret dans le code} (\code{terraform.tfvars}
exclu du dépôt, \code{.example} anonymisé versionné) ; \textbf{state protégé} (\code{gitignore} ; il contient
\code{admin\_password}, \code{primary\_access\_key}, \code{secret} en clair) ; \textbf{Storage privé}, TLS 1.2,
HTTPS-only ; captures \textbf{anonymisées} (IP admin et identifiants Azure floutés). \textbf{Recommandations
production} : Azure Bastion, backend de state distant chiffré + RBAC, \code{allow\_nested\_items\_to\_be\_public
= false}, Application Gateway + WAF en HTTPS, Azure Policy d'obligation de tags.

\subsection{FinOps}
Coût $\approx$ \textbf{46,5 \$/mois en 24/7} (prix réels Azure Retail Prices). Postes principaux : \textbf{Load
Balancer ($\approx$ 18 \$)}, \textbf{compute des 2 VM ($\approx$ 14 \$)}, \textbf{IP publiques ($\approx$
11 \$)} — le LB et les IP facturent même à trafic nul. Leviers : désallouer les VM hors usage, réduire les IP
publiques (Bastion), lifecycle policy sur le Blob, et \code{terraform destroy} entre les séances pour ramener
le coût à $\approx$ 0.

\subsection{Gouvernance et cycle de vie}
\textbf{Workflow} : \code{init} \arr{} \code{fmt} \arr{} \code{validate} \arr{} \code{plan} \arr{} \code{apply}
\arr{} \code{destroy}, le \code{plan} étant systématiquement relu. La modification d'un tag se propage en
\textbf{mise à jour en place} (\code{\textasciitilde}), sans recréation. Une modification manuelle (portail) est
\textbf{détectée} par \code{terraform plan} et \textbf{réconciliée}, démontrant le rôle du code comme
\textbf{source de vérité}.

\subsection{Amélioration proposée (mise en autonomie — Option B : Azure Bastion)}
\textbf{Objectif :} supprimer les IP publiques directes des VM et administrer via \textbf{Azure Bastion}, pour
réduire la surface d'attaque et le coût des IP.
\begin{itemize}
  \item \textbf{Modification :} ajout d'un subnet \code{AzureBastionSubnet} (\code{10.20.255.0/27}) + un \code{azurerm\_bastion\_host} (+ IP publique Standard) ; \textbf{suppression des IP publiques des VM} ; retrait de la règle SSH publique du NSG.
  \item \textbf{Fichiers impactés :} \code{compute.tf} (NIC sans IP publique), \code{network.tf} (subnet Bastion), \code{security.tf} (retrait \code{Allow-SSH-Admin}), nouveau \code{bastion.tf}, \code{variables.tf} (\code{enable\_public\_ip}).
  \item \textbf{Risques :} le Bastion a un \textbf{coût horaire} (à arbitrer en dev) ; le test web direct par VM n'est plus possible (uniquement via le LB) ; la bascule nécessite un \code{plan} attentif (mise à jour des NIC, pas recréation des VM).
\end{itemize}

\subsection*{Conclusion}
Le TP2 fait passer ShopEasy d'une infrastructure \textbf{créée à la main} à une infrastructure \textbf{décrite
en code}, reproductible et gouvernable. Les principes clés — \textbf{moindre privilège}, \textbf{pas de secret
versionné}, \textbf{nommage et tags cohérents}, \textbf{lecture systématique du plan}, \textbf{maîtrise du
drift et des coûts} — sont appliqués et documentés. Les évolutions proposées (Bastion, state distant,
modularisation, CI/CD) rapprocheraient l'environnement d'un usage de production.


% #####################################################################
%  QUIZ
% #####################################################################
\section{Quiz de validation}

\begin{enumerate}[leftmargin=2.2em]
\item \textbf{Terraform est-il un outil impératif ou déclaratif ? Justifier.}\\ \textbf{Déclaratif} : on décrit l'état final souhaité, Terraform calcule les actions (créer/modifier/détruire) pour l'atteindre — d'où l'importance du \code{plan}.
\item \textbf{Quel est le rôle d'un provider Terraform ?}\\ Un plugin qui dialogue avec l'API d'une plateforme et traduit les ressources Terraform en appels concrets (ici \code{azurerm} pour Azure).
\item \textbf{À quoi sert le fichier \code{terraform.tfstate} ?}\\ C'est le fichier d'état : il fait le lien entre le code et les ressources réelles créées dans Azure.
\item \textbf{Pourquoi exécuter \code{terraform plan} avant \code{terraform apply} ?}\\ Pour prévisualiser les changements avant d'impacter Azure ; une destruction ou un remplacement imprévu est un signal d'alerte.
\item \textbf{Quelle commande formate le code Terraform ?}\\ \code{terraform fmt}.
\item \textbf{Quelle commande vérifie la syntaxe et la cohérence de base ?}\\ \code{terraform validate}.
\item \textbf{Quel service Azure correspond au réseau privé logique ?}\\ Le Virtual Network (VNet).
\item \textbf{Quel composant Azure filtre les flux réseau entrants et sortants ?}\\ Le Network Security Group (NSG).
\item \textbf{Pourquoi restreindre SSH à une seule adresse IP ?}\\ Pour réduire la surface d'attaque (scans et force brute) : seule l'IP admin peut tenter une connexion.
\item \textbf{Quel est l'intérêt de \code{count} dans le déploiement des VM ?}\\ Déployer plusieurs instances identiques sans dupliquer le code ; \code{count.index} nomme et relie chaque VM à son interface et son IP.
\item \textbf{Pourquoi utiliser des variables ?}\\ Pour éviter les valeurs codées en dur et rendre le projet paramétrable et réutilisable sur plusieurs environnements.
\item \textbf{Pourquoi utiliser des outputs ?}\\ Pour exposer les informations utiles après déploiement (IP du LB, IP des VM, nom du Storage) sans fouiller le portail.
\item \textbf{Pourquoi taguer les ressources ?}\\ Pour la gouvernance et le FinOps : imputer les coûts, filtrer les ressources, identifier le propriétaire et l'environnement.
\item \textbf{Différence entre un changement Terraform et un changement manuel dans le portail ?}\\ Le changement Terraform est versionné, relu et tracé (\code{plan}/\code{apply}) ; le changement manuel se fait hors code \arr{} dérive non tracée.
\item \textbf{Quel risque présente un Storage Account public ?}\\ Une fuite de données : les blobs deviennent accessibles à quiconque connaît l'URL (risque RGPD).
\item \textbf{Pourquoi le versioning Blob peut-il générer des coûts supplémentaires ?}\\ Il conserve une copie de chaque version à chaque modification ; ces versions s'accumulent et sont facturées comme du stockage.
\item \textbf{À quoi sert un Load Balancer ?}\\ À répartir le trafic entrant sur plusieurs VM derrière un point d'entrée unique, ce qui améliore la disponibilité et permet de monter en charge.
\item \textbf{Pourquoi un state distant est-il préférable en équipe ?}\\ Emplacement central partagé, verrouillage (pas de corruption), pas de state sur les postes, intégration CI/CD.
\item \textbf{Citer deux bonnes pratiques de sécurité pour un projet Terraform.}\\ Ne jamais versionner de secret (\code{tfvars}/\code{tfstate} exclus, Key Vault) ; moindre privilège (SSH restreint, NSG, backend de state distant sécurisé).
\item \textbf{Citer deux améliorations pour rapprocher l'architecture d'un environnement de production.}\\ Supprimer les IP publiques des VM + Azure Bastion et externaliser le state ; Application Gateway + WAF (HTTPS), supervision (Azure Monitor) et pipeline CI/CD.
\end{enumerate}

\vspace{6pt}
\begin{tcolorbox}[enhanced, sharp corners, boxrule=0pt, leftrule=4pt, colback=okgreenbg, colframe=okgreen,
   left=4.5mm,right=4.5mm,top=3.6mm,bottom=3.6mm]
{\sffamily\bfseries\color{okgreen}\eParty~Fin du dossier} — les 14 ateliers, la note technique et le quiz de
validation sont complets : projet Terraform initialisé et paramétré, déploiement Azure (réseau, NSG, 2 VM
cloud-init, Load Balancer, Storage privé versionné), outputs, modification et dérive maîtrisées, state distant,
analyses coût/sécurité/maintenabilité et nettoyage.
\end{tcolorbox}

\end{document}
